\section{\texorpdfstring{$X_{12}$}{X12} の Poisson 概型の分類}
\label{p12:classification}
\subsection{設定、主結果と先行研究}
$X=\PP^1\times\PP^2$ とする。
Polishchuk~\cite[定理6.1, 式(16)]{Polishchuk1997} は、階数2束の行列式の Poisson 加群構造を固定すると、
その射影束上で基底の双ベクトルと両立する Poisson 構造を、束の Poisson 加群構造で記述できることを示した。
本稿では $Y=\PP^2$、束 $\mathcal O_Y^{\oplus2}$ にこの対応を適用する。
平面の退化三次曲線の補集合で得られる平坦 $\mathrm{SL}_2$ 接続の局所モノドロミーと、
混合項の次数制約を組み合わせることが、各ファイバーの分類の要点である。
Loray--Pereira--Touzet~\cite[定理9.1]{p12-LPT} は Picard 数1の Fano 三次元多様体上の Poisson 構造を分類した。
Belmans--Fatighenti--Tanturri~\cite[\S2.2]{p12-BFT} は高い Picard 数での分類問題と多重ベクトル場の次元を論じている。
本節で扱う $X$ の Picard 数は2である。
Pym~\cite{p12-Pym} は1形式や葉層を用いる手法を概説している。
Okumura~\cite[定理1.2]{p12-Okumura} の積に関する結果は偶数次元で各因子の次元が3以上の場合の
SNC log symplectic 構造を扱う。以下の個々の成分について既存文献との優先関係は未確定である。

接束の分解と $\bigwedge^2T_{\PP^2}\cong\mathcal O_{\PP^2}(3)$ より
\[
 H^0(X,\textstyle\bigwedge^2T_X)
 \cong (\mathfrak{sl}_2\otimes\mathfrak{sl}_3)\oplus H^0(\PP^2,\mathcal O(3)),
 \qquad \dim=24+10=34.
\]
また $-K_X=\mathcal O(2,3)$ より $h^0(X,\bigwedge^3T_X)=30$ である。

\begin{definition}\label{p12:def:scheme}
$V=H^0(X,\bigwedge^2T_X)$ と $W=H^0(X,\bigwedge^3T_X)$ を置く。
本節の Poisson 概型は、二次写像
\[
 V\longrightarrow W,\qquad \pi\longmapsto\tfrac12[\pi,\pi]
\]
の零ファイバー $\operatorname{Pois}(X)$ である。
したがって、その座標環は $\operatorname{Sym}(V^\vee)$ を
$\tfrac12[\pi,\pi]$ の全座標二次式で割った環である。
既約成分 $Z$ の生成点 $\eta$ における重複度は
$\operatorname{length}\mathcal O_{\operatorname{Pois}(X),\eta}$ と定める。
以下の分類では固定した $X$ 上の双ベクトルそのものをパラメータとする。
\end{definition}


本節で分類する成分の一覧を掲げる。次元は
$\operatorname{Aut}(X)$ で割る前のアフィン Poisson 概型の次元を表す。
\[
\begin{array}{c|c|l}
\text{記号}&\text{次元}&\text{構成または特徴}\\\hline
V_{\mathrm{vert}}&10&\PP^2\text{ 上の Poisson 構造の引き戻し}\\
C_{\mathrm{tor}}&11&\text{最大トーラスの双ベクトル}\\
C_{\mathrm{fib}}&11&H\omega,\ H\in H^0(\mathcal O(2,1)),\ \omega\in H^0(\Omega^1_{\PP^2}(2))\\
C_{\mathrm{line\text{-}conic}}&11&aLQ\,du+h(u)(L\,dQ-2Q\,dL)\\
C_{(1,1)^2(0,1)}&11&\text{$(1,1)$ 因子を葉とする pencil}\\
C_{(1,1)(1,0)(0,2)}&12&\text{$(1,1)$ 因子を葉とする logarithmic 型}\\
C_{\mathrm{Ric}}&12&\text{共点三直線と直線束上の Riccati 方程式}\\
C_4&12&\text{四因子の対数形式、二つの自由な留数}\\
C_{(1,2)(1,0)(0,1)}&14&\text{$(1,2)$ 因子を葉とする logarithmic 型}
\end{array}
\]
\begin{theorem}\label{p12:th:complete}
$\mathbb C$ 上で、上表の九つが $\operatorname{Pois}(X)$ の全既約成分である。
$C_{\mathrm{Ric}}$ の生成点での長さは2であり、他の八つの生成点での長さは1である。
\end{theorem}
網羅性は射影三次式による層別化から従う。証明を第\ref{p12:sec:conclusion}節でまとめる。
成分の一般点の重複度と、二つの指定因子への制限ファイバーの概型構造は別の問題であり、
後者は定理\ref{p12:th:extension5}, \ref{p12:th:extensionfibers} と命題\ref{p12:prop:zero-restriction} で扱う。

\subsection{アフィン表示と Jacobi 方程式}
標準アフィン開集合の座標を $u$（$\PP^1$）および $v,w$（$\PP^2$）とし、
\[
 \pi=A\,\partial_u\wedge\partial_v+B\,\partial_u\wedge\partial_w+C\,\partial_v\wedge\partial_w.
\]
ここで $C=C(v,w)$ は次数 $\le3$ の多項式であり、$(A,B)$ は $\mathfrak{sl}_2\otimes\mathfrak{sl}_3$ の元を表す。$\mathfrak{sl}_2$ 側の基底は $1,u,u^2$、$\mathfrak{sl}_3$ 側の8つのベクトル場の基底は
\[
 \partial_v,\quad\partial_w,\quad v\partial_v,\quad w\partial_v,\quad v\partial_w,\quad w\partial_w,
 \quad v^2\partial_v+vw\partial_w,\quad vw\partial_v+w^2\partial_w.
\]
Schouten 条件 $[\pi,\pi]=0$ は、Jacobi 恒等式により次の多項式の恒等的消滅と同値である。
\begin{equation}\label{p12:eq:jacobi}
 J(A,B,C):=A(C_v+B_u)+B(C_w-A_u)-C(B_w+A_v)=0.
\end{equation}
$J$ は $u$ に関して次数 $\le2$、$v,w$ に関して全次数 $\le3$ であり、係数比較によって30個の斉次二次式を与える。

\subsection{垂直成分}
\begin{theorem}\label{p12:th:vertical}
 $V_{\mathrm{vert}}=\{(A,B,C):A=B=0\}\cong\mathbb{A}^{10}$ は $\operatorname{Pois}(X)$ の既約成分であり、滑らかな平面三次曲線を表す一般点で $\operatorname{Pois}(X)$ は滑らかである。
\end{theorem}
\begin{proof}
 $A=B=0$ なら \eqref{p12:eq:jacobi} は恒等的に零である。$C$ が滑らかな平面三次曲線を定める点での一次変形を $\dot\pi=\dot C+\eta$ ($\eta\in\mathfrak{sl}_2\otimes\mathfrak{sl}_3$) とする。一次 Jacobi 条件は
 $[C,\eta]=0$ である。$\eta=\sum_i a_i\wedge b_i$ と表せば、これは
 \[
 \sum_i a_i\otimes\mathcal L_{b_i} C=0
 \]
 となる（符号規約は零条件に影響しない）。滑らかな平面三次曲線の射影自己同型群は有限であるから、$b\mapsto\mathcal L_bC$ は $\mathfrak{sl}_3\to H^0(\PP^2,\mathcal O(3))$ の単射である。したがって $\eta=0$。接空間は $\dot C$ の10次元空間に一致する。$V_{\mathrm{vert}}$ 自身の次元も10なので、この点は $\operatorname{Pois}(X)$ の滑らかな点であり、$V_{\mathrm{vert}}$ は既約成分である。
\end{proof}

\subsection{トーラス由来の成分}
$E_u=u\partial_u$, $E_v=v\partial_v$, $E_w=w\partial_w$ は $\mathrm{Aut}(X)$ の最大トーラスの大域ベクトル場である。
\[
 \pi_{\lambda,\mu,\nu}=\lambda E_u\wedge E_v+\mu E_u\wedge E_w+\nu E_v\wedge E_w
\]
は任意の $\lambda,\mu,\nu\in\mathbb C$ について Poisson 構造をなす。$G=\mathrm{PGL}_2\times\mathrm{PGL}_3$ とし、$C_{\mathrm{tor}}$ を
\[
 G\cdot\{\pi_{\lambda,\mu,\nu}\mid(\lambda,\mu,\nu)\in\mathbb C^3\}
\]
の Zariski 閉包と定める。
\begin{theorem}\label{p12:th:toric}
 $C_{\mathrm{tor}}$ は $\operatorname{Pois}(X)$ の11次元既約成分である。$C_{\mathrm{tor}}\ne V_{\mathrm{vert}}$ である。
\end{theorem}
\begin{proof}
 $G\times\mathbb A^3$ の既約性から $C_{\mathrm{tor}}$ は既約である。点
 \[
 \pi_0=2uv\,\partial_u\wedge\partial_v+3uw\,\partial_u\wedge\partial_w+5vw\,\partial_v\wedge\partial_w
 \]
 を考える。$G$ の Lie 代数のうち最大トーラスに属する3次元部分は $\pi_0$ を固定する。残りの root direction は $\mathfrak{sl}_2$ で2つ、$\mathfrak{sl}_3$ で6つあり、$\mathcal L_V\pi_0$ はそれぞれ互いに異なる非零トーラス重みを持つ（これらが非零であることは上の座標表示への代入で直接確認できる）。したがって軌道の接空間は8次元である。$\lambda,\mu,\nu$ に関する3つの変分は互いに独立な重み0のベクトルなので、$C_{\mathrm{tor}}$ の $\pi_0$ における局所次元は少なくとも11である。

 一方、\eqref{p12:eq:jacobi} を $\pi_0$ で線形化する。$\dot\pi$ の30係数への線形写像の rank が23以上であることを示す。トーラス重みによって各項を分解する。$\dot C=v^i w^j$（$i+j\le3$）とすると、線形化された Jacobi 多項式は
 \[
 \dot J=u\bigl(2(i-1)+3(j-1)\bigr)v^i w^j.
 \]
 唯一 $(i,j)=(1,1)$ のとき係数が消え、残り9つは独立な重みの非零像となる。

 次に $\dot C=0$ とし、$(\dot A,\dot B)=u^k(P,Q)$（$k=0,2$）で $P\partial_v+Q\partial_w$ を $\mathfrak{sl}_3$ の前節の基底に走らせる。$k=0$ では像は順に
 \[
 8w,\ 3v,\ 3vw,\ 8w^2,\ 3v^2,\ -2vw,\ -4v^2w,\ -4vw^2,
 \]
 $k=2$ では順に
 \[
 u^2\bigl(2w,\ 7v,\ -3vw,\ 2w^2,\ 7v^2,\ 2vw,\ -6v^2w,\ -6vw^2\bigr).
 \]
 各 $k$ について7つの異なる単項式が得られ、合計14次元の像を作る。先の9つの像は $u$ 次数が1であるため、これらと独立である。ゆえに rank $\ge9+14=23$、すなわち
 \[
 \dim T_{\pi_0}\operatorname{Pois}(X)\le34-23=11.
 \]
 これと11次元の族の存在より $\dim T_{\pi_0}\operatorname{Pois}(X)=11$、かつ $\pi_0$ は滑らかな点である。よってその点を通る11次元既約閉集合 $C_{\mathrm{tor}}$ は既約成分である。$\pi_0$ は mixed term を持つので $V_{\mathrm{vert}}$ とは異なる。
\end{proof}


\subsection{既約三次曲線上の剛性と零ファイバー}
$Y=\PP^2$ とし、射影に伴う線形写像
$q:\operatorname{Pois}(X)\to H^0(Y,-K_Y)$ で平面成分を取り出す。
$D=(C=0)$ が被約既約な三次曲線である $C$ の全体を $V^{\mathrm{irr}}$ と書く。
これは原点と直線・二次式の積からなる閉集合の補集合である。

\begin{theorem}\label{p12:th:irreducible-cubic}
$q^{-1}(V^{\mathrm{irr}})\to V^{\mathrm{irr}}$ は概型の同型であり、
その逆は混合項を零とする切断である。特に $D$ が滑らかな場合のほか、
既約な結節三次曲線および尖点三次曲線の場合にも非零混合項は存在しない。
\end{theorem}
\begin{proof}
$C$ に対応する平面の双ベクトルを $\sigma$、その非退化開集合を
$U=Y\setminus D$ とする。$\pi_1(U)=\mathbb Z/3$ は三つの型すべてについて成り立つ
\cite[\S2.3.2, (3e)]{p12-Cogolludo}。滑らかな場合は
$\{C=1\}\to U$ が $\mu_3$ 被覆で、斉次孤立特異点の Milnor
ファイバーが単連結であることからも従う~\cite[\S\S5--6]{p12-Milnor}。

$\pi$ の混合項を $\sum_{i=1}^3 e_i\wedge V_i$ と書く。
ここで $e_i$ は $H^0(\PP^1,T_{\PP^1})$ の基底、$V_i$ は $Y$ の射影ベクトル場である。
$\PP^1\times U$ 上の特性分布は $TU$ に同型に射影するので、
自明な $\PP^1$ 束の接続を定める。水平持ち上げの垂直成分は、符号の取り方を除いて
$\sum_i (\iota_{V_i}\sigma^{-1})\,e_i$ である。Jacobi 恒等式により水平分布は可積分である。
射影ベクトル場の Lie 環 $\mathfrak{pgl}_2$ を $\mathfrak{sl}_2$ と同一視すると、
これは自明な階数2束上の平坦な接続 $d+\Omega$ を与える。
すなわち $\Omega$ を水平持ち上げから符号を合わせて取ると、曲率の射影像が零であり、
曲率はトレース零だから曲率自身も零である。
この対応は Poisson 加群の射影化~\cite{Polishchuk1997} の局所表示でもある。

$D$ の滑らかな点の近傍で $D=(x=0)$ となる座標 $(x,y)$ を取る。
$\sigma^{-1}$ は1位の極しか持たないので
$\Omega=(A\,dx+B\,dy)/x$ と書ける。$A,B$ は正則なトレース零行列である。
平坦性の $x^{-2}$ の係数は、$A_0=A(0,y),\,B_0=B(0,y)$ として
\[
 [A_0,B_0]=B_0
\]
となる。ある点で $B_0\ne0$ なら $\operatorname{ad} A_0$ が固有値1を持ち、
$A_0$ の二つの固有値は $1/2,-1/2$ である。この点で $y$ を固定し、
$x=r e^{i\theta}$ に沿った平行移動を考える。その行列の微分方程式の係数は
$-i A(r e^{i\theta},y)$ であり、$r\to0$ で一様に $-i A_0$ に収束する。
円周のホロノミー行列も $\exp(-2\pi i A_0)$ に収束する。
各ホロノミー行列は3乗が単位行列であるが、極限の固有値は $-1$ となり矛盾する。
したがって $B_0=0$ で、$\Omega$ は $D$ の滑らかな点で対数極を持つ。
$\sigma=x\,h\partial_x\wedge\partial_y$ ($h$ は単元) と書けば、
$\iota_{V_i}\sigma^{-1}$ の $dy$ 成分が正則であることは
$V_i(x)\in(x)$ と同値である。よって $V_i$ は $D$ を保存する。

$\mathcal L_V\sigma=0$ を満たす射影ベクトル場は零である
\cite[Prop.~15.1]{Polishchuk1997}。実際、$V$ は $D$ に接し、
$\alpha=\iota_V\sigma^{-1}$ は $U$ 上で閉じており、$D$ の滑らかな点で
対数極を持つ。$\pi_1(U)$ が有限なので任意の閉路に沿う周期は零である。
とくに局所メリディアンの周期が零だから対数留数が消え、$\alpha$ は正則に延びる。
残る有限個の特異点にも Hartogs の定理で延び、
$H^0(Y,\Omega_Y^1)=0$ より $V=0$ となる。
一方、$D$ を保存するベクトル場 $V$ は $\mathcal L_V\sigma=c_V\sigma$
を満たす。実際、左辺は $D$ 上消える $-K_Y$ の切断である。
いま示したことから $V\mapsto c_V\in\mathbb C$ は単射なので、
$D$ を保存する射影ベクトル場の空間は高々1次元である。
従って混合項が非零なら $e\wedge V$ と書けるが、このとき
式~\eqref{p12:eq:jacobi} は $\mathcal L_V\sigma=0$ を与え、矛盾する。

これで閉点の主張を得た。各閉点 $(0,C)$ での一次 Jacobi 条件は
各 $i$ について $\mathcal L_{\dot V_i}\sigma=0$ であり、混合方向は零である。
垂直な10次元部分空間が通り、接空間も10次元なので局所環は正則である。
従ってこの開部分概型は被約で、上の切断による同定は概型として成立する。
\end{proof}

\begin{proposition}\label{p12:prop:zero-cubic}
$q^{-1}(0)_{\mathrm{red}}$ の既約成分は次の二つである。
第一はテンソル階数が高々1の10次元錐であり $C_{\mathrm{tor}}$ に含まれる。
第二は $p\in\PP^2$ を動かして、$p$ を通る直線束に接する三つの射影ベクトル場を
任意に選ぶ11次元族であり、$C_{\mathrm{fib}}$ に一致する。
\end{proposition}
\begin{proof}
$\partial_u\wedge (V_0+uV_1+u^2 V_2)$ の Jacobi 条件は
$(V_0+uV_1+u^2 V_2)\wedge(V_1+2uV_2)=0$ であり、
各 $V_i\wedge V_j=0$ と同値である。階数が2以上なら、
どの非零の $V\in\langle V_i\rangle$ に対しても、これと独立な $W$ が取れ、
$W=fV$ を満たす有理関数 $f$ がある。$V$ に因子的な零点がなければ
$f$ は余次元1の全点で正則、従って定数であり矛盾する。
行列 $M$ が定める射影ベクトル場の零点は固有空間の射影化の和である。
直線上で消える非零射影ベクトル場には、スカラー行列を引いて
階数1の表示 $M=p\otimes\ell$ がある。その一般点での方向は
その点と $[p]$ を結ぶ直線の方向であるから、すべての $V_i$ が同じ中心を持つ。
逆に、中心 $[p]$ に対し
$W_p=\{p\otimes\ell\bmod\mathbb C I:\ell\in(\mathbb C^3)^*\}$
内の任意の三つ組は条件を満たす。
この族は $\PP^2$ 上の階数9ベクトル部分束の全空間からの射影像なので
既約閉であり、階数2以上では中心は一意なので次元は11である。
中心を固定した直線束の1形式に $H\in H^0(\mathcal O(2,1))$ を掛けたものが
ちょうどこの族である。

第一の族は Segre 錐なので次元 $3+8-1=10$ である。その一般点では
両因子のベクトル場が半単純で、最大トーラスに属するから閉包全体が
$C_{\mathrm{tor}}$ に入る。一般の $\mathfrak{sl}_3$ の元は相異なる三つの
固有値を持ち、因子的な零点を持たないので第二の族には入らない。
第二の族には階数3のテンソルが存在する。この二つは互いに包含されず、
零ファイバーの既約成分を尽くす。
\end{proof}


\subsection{二つの del Pezzo 因子の幾何}
$H_1,H_2$ を $X$ の二つの因子から引いた超平面類とする。
$-K_X=2H_1+3H_2$ であり、$(1,a)$ 型の滑らかな因子 $S$ に対し
$-K_S=(H_1+(3-a)H_2)|_S$ が成り立つ。その次数は
\[
((-K_X-S)^2\cdot S)_X=
\begin{cases}8&(a=1),\\5&(a=2)\end{cases}
\]
と直接交点積を取ればよい。
射影 $S\to\PP^2$ は、二つの $a$ 次形式の基点を爆発したものだから、
一般の $a=1$ では1点を爆発して $S\cong\mathbb F_1$（8次）、
$a=2$ では一般位置の4点を爆発して5次 del Pezzo となる。
それぞれ
\[
H^0(S,-K_S)\cong
\begin{cases}
H^0(S,\mathcal O_S(1,2)),\quad h^0=9 &(a=1),\\
H^0(S,\mathcal O_S(1,1)),\quad h^0=6 &(a=2).
\end{cases}
\]
曲面では正則双ベクトルはすべて Poisson である。

\subsubsection{Poisson 因子の延長問題}
滑らかな因子 $S\subset X$ を Poisson 部分多様体とするための条件は、
イデアル層 $I_S$ が Poisson イデアルとなること、すなわち
\[
\{I_S,\mathcal O_X\}\subset I_S
\]
である\cite[\S4.1]{p12-Pym}。これは $\pi|_S\in\bigwedge^2T_S$ と同値である。
この条件を満たす双ベクトルの層を $\bigwedge^2T_X(-\log S)$ と表すと、
\begin{equation}\label{p12:eq:log-exact}
0\longrightarrow \bigwedge^2T_X(-S)
\longrightarrow \bigwedge^2T_X(-\log S)
\longrightarrow \bigwedge^2T_S\longrightarrow0
\end{equation}
が成り立つ（局所座標 $S=\{z=0\}$ で直接検証できる）。

\begin{theorem}\label{p12:th:extension5}
$S$ を滑らかな5次 del Pezzo 型の $(1,2)$ 因子とする。
任意の $\beta\in H^0(S,-K_S)$ は、$S$ に接する大域双ベクトル
$\widetilde\beta\in H^0(X,\bigwedge^2T_X)$ に\emph{一意に}延長する。
ただし $\widetilde\beta$ 自体が Poisson になる必要十分条件は
$\beta$ に対応する $H\in H^0(X,\mathcal O_X(1,1))\cong\mathbb C^2\otimes\mathbb C^3$
のテンソル階数が1以下であることである。
従って、延長可能な $\beta$ の集合は4次元の Segre 錐である。
\end{theorem}
\begin{proof}
接束の分解を用いると、$S\sim H_1+2H_2$ に対し
\[
\bigwedge^2T_X(-S)
=\mathcal O(-1,1)\oplus
\bigl(\mathcal O_{\PP^1}(1)\boxtimes T_{\PP^2}(-2)\bigr).
\]
$H^0(\PP^2,T_{\PP^2}(-2))=H^1(\PP^2,T_{\PP^2}(-2))=0$ は
捻った Euler 列から従う。$\PP^1$ 上の $\mathcal O(-1)$ も $H^0,H^1$ が零なので、
$H^0,H^1$ ともに $\bigwedge^2T_X(-S)$ では零になる。
よって \eqref{p12:eq:log-exact} の大域切断により
\[
H^0(X,\bigwedge^2T_X(-\log S))\ \xrightarrow{\sim}\ H^0(S,-K_S).
\]
一方、完全列 $0\to\mathcal O_X(0,-1)\to\mathcal O_X(1,1)\to\mathcal O_S(1,1)\to0$
から、$H^0(X,\mathcal O_X(1,1))\xrightarrow{\sim}H^0(S,-K_S)$ も従う。

$F\in H^0(X,\mathcal O(1,2))$ を $S$ の定義式、
$H\in H^0(X,\mathcal O(1,1))$ を $\beta$ に対応する切断とする。
積多重次数を持つアフィン錐上で、二つの射影因子の方向の外微分を $d_1,d_2$ と書くと、
\begin{equation}\label{p12:eq:formdp5}
\alpha(F,H)=H(d_1F+d_2F)-F(d_1H+2d_2H)
\end{equation}
は両方の Euler ベクトル場との縮約が零なので、
$\Omega^1_X(2,3)\cong\bigwedge^2T_X$ の大域切断に降下する。
$S$ への制限では $\alpha|_S=H\,dF$ であり、必要な $\widetilde\beta$ に対応する。

$S$ が5次 del Pezzo であるとき、四つの一般位置の基点は $\mathrm{PGL}_3$ の作用で標準形へ移せる。
従って $X$ の自己同型により、アフィン座標 $(u,v,w)$ で
\[
F=uvw+(v-1)(w-1)
\]
としてよい。$H$ を
\[
H=(a_0+a_1v+a_2w)+u(a_3+a_4v+a_5w)
\]
とし、$\alpha=C\,du-B\,dv+A\,dw$ と置く。
\eqref{p12:eq:jacobi} に代入して整理すると、
\begin{align}\label{p12:eq:minoridentity}
J(A,B,C)=F\bigl[&M_{01}((u+1)v-1)
-M_{02}((u+1)w-1)\\
&-M_{12}(v+w-2)\bigr],\nonumber
\end{align}
ただし
\[
M_{01}=a_0a_4-a_1a_3,\quad
M_{02}=a_0a_5-a_2a_3,\quad
M_{12}=a_1a_5-a_2a_4.
\]
従って $J=0$ は三つの $2\times2$ 小行列式 $M_{01}=M_{02}=M_{12}=0$
と同値である。このアフィン代数集合は $2\times3$ 行列の階数1以下の錐、
すなわち $\operatorname{Cone}(\PP^1\times\PP^2)$（次元4）である。
式 \eqref{p12:eq:minoridentity} は付属コードでも記号的に検証している。
\end{proof}

\begin{remark}
上の $H=U\cdot L$ が分解可能な場合、\eqref{p12:eq:formdp5} は
\[
\alpha=UL\,dF-FL\,dU-2FU\,dL
\]
となる。これは $F/(UL^2)$ を有理第一積分とする Poisson 構造であり、
固定した $S$ に限らず、その pencil に含まれる一般の5次 del Pezzo 因子を Poisson 葉の閉包とする。
\end{remark}

\subsubsection{一般の del Pezzo Poisson 構造の延長障害}
次は5次・8次の両方に有効な制約である。零曲線が特異な既約曲線であっても成立する。
\begin{proposition}\label{p12:prop:flatobstruction}
$S\subset X$ を滑らかな $(1,1)$ または $(1,2)$ 型 del Pezzo 因子とする。
$\beta\in H^0(S,-K_S)$ の零因子 $E=\{\beta=0\}$ が既約なら、
$\beta$ を誘導して $S$ を Poisson 部分多様体とする ambient Poisson 構造は存在しない。
\end{proposition}
\begin{proof}
もし延長 $\pi$ が存在すれば、Poisson イデアル $I_S$ の商
$I_S/I_S^2\cong N_{S/X}^{\vee}$ は $S$ 上の Poisson module になる。
$S\setminus E$ は正則シンプレクティックなので、$N_{S/X}^{\vee}|_{S\setminus E}$
は平坦な代数的接続を持ち、従ってその de Rham 第一 Chern 類は零である。
$E$ が既約なら、局在長完全列と Borel--Moore 双対性により
\[
\ker\bigl(H^2_{\mathrm{dR}}(S)\to H^2_{\mathrm{dR}}(S\setminus E)\bigr)
=\mathbb C[E].
\]
従って $c_1(N_{S/X})$ は $[E]=c_1(-K_S)$ に比例しなければならない。
しかし $N_{S/X}=\mathcal O_S(1,a)$ と $-K_S=\mathcal O_S(1,3-a)$ は
$a=1,2$ のいずれについても独立な Picard 類である。これは矛盾する。
Poisson module と平坦性の標準的背景は\cite{p12-GualtieriPym,p12-Pym} を参照。
\end{proof}


\subsubsection{固定した8次 del Pezzo 因子での完全な延長判定}
$S$ を滑らかな $(1,1)$ 因子とする。自己同型で
\[
 F=uv+w=0
\]
とおける。$p=(v,w)=(0,0)$ とすると
$S=\operatorname{Bl}_{p}\PP^2$ であり、$S$ 上の反標準切断は
$p$ を通る平面三次曲線の方程式 $P(v,w)$ に対応する。

\begin{theorem}\label{p12:th:fixed8}
$S\subset X$ を滑らかな $(1,1)$ 因子とし、
$\beta\in H^0(S,-K_S)$ を非零切断とする。
$\beta$ が $S$ を Poisson 部分多様体とする
$X$ 上の正則 Poisson 双ベクトルへ延長するための必要十分条件は、
対応する平面三次式 $P$ が一次式と二次式の積に分解することである（重複因子を許す）。
従って、延長可能な切断の集合は、$H^0(S,-K_S)\simeq\mathbb{A}^9$ 内の
相異なる二つの7次元既約閉錐
\[
 \mathcal R_{\ell_p q}=\{\ell_p q\mid p\in\ell_p,\deg q=2\},
 \qquad
 \mathcal R_{\ell q_p}=\{\ell q_p\mid p\in q_p,\deg \ell=1\}
\]
の和集合である（$0$ も含む）。
\end{theorem}
\begin{proof}
$S\to\PP^2$ が同型となる稠密開集合で見ると、平面への射影 $q(\pi)$ は、
指定した $\beta$ に対応する三次式 $P$ に一致する（標準化での非零定数倍を許す）。
$P$ が既約なら定理\ref{p12:th:irreducible-cubic} より $\pi$ は平面方向のみである。
$P\ne0$ となる点でその像は平面全体の接方向であり、
非定数有理関数 $u=-w/v$ のグラフ $S$ に接することはできない。
これで必要性が従う。

後で制限ファイバー自体を計算するため、接する一般の双ベクトルを記述する。
$\bigwedge^2T_X\simeq\Omega^1_X(2,3)$ の対応で、
$S$ に接する一般の双ベクトルは以下の1形式で表せる。
\begin{equation}\label{p12:eq:normal8}
 \alpha=2H\,dF-F(2d_1H+d_2H)+F\gamma,
\end{equation}
ただし $d_i$ はそれぞれの射影因子の斉次微分であり、
\begin{align*}
H={}&a_0+a_1u+v(b_0+b_1u)+b_2uw
+v^2(c_0+c_1u)+c_2uvw+c_3uw^2,\\
\gamma={}&(x_0+x_1u)\,dv+(y_0+y_1u)\,dw
+(z_0+z_1u)(v\,dw-w\,dv).
\end{align*}
この表示が全切断を尽くすことを示す。接する双ベクトルの完全列
\eqref{p12:eq:log-exact} の核は
\[
H^0(X,\bigwedge^2T_X(-S))
 \cong H^0(\PP^1,\mathcal O(1))\otimes
 H^0(\PP^2,T_{\PP^2}(-1)),\qquad\dim=6
\]
であり、$H^1$ は消える。
従って大域切断空間の次元は $6+9=15$ である。
\eqref{p12:eq:normal8} は15個の独立な係数を含み、核と反標準切断の
両方を尽くす（付属コードで直接の線形代数検証も行う）。

$\alpha=C\,du-B\,dv+A\,dw$ として
\eqref{p12:eq:jacobi} に代入すると、$J(A,B,C)$ は $F$ で割れ、
\[
 J(A,B,C)=F\sum_{i,j,k}p_{ijk}u^iv^jw^k,
\]
ただし指数は $0\le i\le1$, $j,k\ge0$, $j+k\le2$ を満たす12組にわたる。
$\alpha|_S=2H\,dF$ であり、平面三次曲線のアフィン方程式は
\begin{align}\label{p12:eq:cubic8}
P(v,w)={}&a_0v-a_1w+b_0v^2-b_1vw-b_2w^2\nonumber\\
 &+c_0v^3-c_1v^2w-c_2vw^2+c_3w^3.
\end{align}
逆に $C$ が可約なら $C=\ell q$ と因数分解し、$p\in C$ から
$p\in\ell$ または $p\in q$ である。
前者では $\ell$ に対応する $U\in H^0(\mathcal O_X(1,0))$ と
$q$ に対応する $Q\in H^0(\mathcal O_X(0,2))$ を取り、
\[
 \alpha=2UQ\,dF-2FQ\,dU-FU\,dQ
\]
を用いる。後者では $q$ の狭義変換が
$\mathcal O_S(1,1)$ の切断に対応するので
$G\in H^0(\mathcal O_X(1,1))$ と $L\in H^0(\mathcal O_X(0,1))$
を取り、
\[
 \alpha=L(G\,dF-F\,dG)
\]
を用いる。いずれも可積分で $S$ に所望の Poisson 構造を誘導する
（スカラー倍は調整できる）。
最後に、第一の可約三次曲線の射影パラメータ空間は
$\PP^1\times\PP^5$、第二は $\PP^2\times\PP^4$ なので、
両錐はそれぞれ7次元であり、一般点の因数分解の型が違うから相異なる。
\end{proof}

\begin{proposition}\label{p12:prop:zero-restriction}
固定した滑らかな $(1,1)$ 因子 $S$ に Poisson 因子として接し、
$S$ 上では零 Poisson 構造を誘導する ambient Poisson 構造は、
\eqref{p12:eq:normal8} で $H=0$ とした $\alpha=F\gamma$ のうち
\[
\operatorname{rank}\begin{pmatrix}x_0&y_0&z_0\\x_1&y_1&z_1\end{pmatrix}\le1
\]
を満たすものに正確に一致する。これは6次元線形空間中の4次元 Segre 錐である。
\end{proposition}
\begin{proof}
$H=0$ なら $\alpha=F\gamma$ であり、
$\alpha\wedge d\alpha=F^2\gamma\wedge d\gamma$。
$\gamma=\gamma_0+u\gamma_1$、
$\gamma_i\in H^0(\PP^2,\Omega^1(2))$ と書くと、
$\gamma\wedge d\gamma$ が零であることは
$\gamma_0\wedge\gamma_1=0$ に等しい。
自然な写像
$\bigwedge^2H^0(\PP^2,\Omega^1(2))\longrightarrow
H^0(\PP^2,\Omega^2(4))$ は3次元空間の同型であるため、
$\gamma_0,\gamma_1$ は比例する。これは表示した階数条件と同値である。
\end{proof}

\subsubsection{8次 del Pezzo 因子での延長の個数とスキーム構造}
\label{p12:sec:markedfibers}
滑らかな $(1,1)$ 因子 $S=\operatorname{Bl}_p\PP^2$ を固定する。
$\mathcal E_\beta$ は、$S$ に接し、指定した切断
$\beta\in H^0(S,-K_S)$ に制限される大域双ベクトルのアフィン空間の中で、
Jacobi 方程式の全係数を課して得られる閉部分スキームとする。
スカラー倍も自己同型も同一視しない。
以下、$D_2=\operatorname{Spec}\mathbb C[\epsilon]/(\epsilon^2)$、
$N=\operatorname{Spec}\mathbb C[s,t]/(st)$、$\mathrm{pt}=\operatorname{Spec}\mathbb C$ とする。

\begin{theorem}\label{p12:th:extensionfibers}
非零 $\beta$ に対応する三次曲線 $C$ が可約のとき、
$\mathcal E_\beta$ のスキーム同型型は以下の通りである。
$L$ は直線、$Q$ は滑らかな二次曲線、$M$ は $L$ と異なる直線を表す。
\[
\begin{array}{lll}
C & p\text{ の位置} & \mathcal E_\beta\\ \hline
L+Q & \text{一方の成分だけに属する} & \mathrm{pt}\\
L+Q & L\cap Q & \mathrm{pt}\sqcup\mathrm{pt}\\
\text{三角形} & \text{辺の滑らかな点} & \mathbb A^1\\
\text{三角形} & \text{頂点} & N\\
\text{共点の相異なる三直線} & \text{共点以外} & \mathbb A^1\\
\text{共点の相異なる三直線} & \text{共点} & \mathbb A^2\\
2L+M & L\setminus M & \mathbb A^1\times D_2\\
2L+M & M\setminus L & \mathbb A^2\\
2L+M & L\cap M & \mathbb A^1\times N\\
3L & L & Z_3
\end{array}
\]
ここで $Z_3=\operatorname{Spec}\mathbb C[A,B,C',D]/(AD-2BC',B^2,BD)$ である。
最初の二行は $L$ が割線の場合にも接線の場合にも成立する。
特に定理\ref{p12:th:fixed8} の二つの7次元錐のいずれでも一般の切断の
延長は一意で、そのファイバーは被約である。最後の行の台は二つの平面の和であり、
$D=0$ の平面の一般点で重複度2を持つ。
\end{theorem}

\begin{proof}
$S\to\PP^1$ は $p$ を中心とする直線束への射影であるため、
標識点を保つ平面の射影変換はその直線束上の変換と合わせて $(X,S)$ の
自己同型に持ち上がる。さらに $\beta$ を非零定数倍すると、
すべての延長を同じ定数倍するスキーム同型が得られる。
従って曲線と点の対 $(C,p)$ の射影同値類ごとに計算すればよい。

$L+Q$ の場合は $L$ が割線か接線かを区別する。
二次曲線の自己同型群 $\operatorname{PGL}_2$ の中で、
割線の安定化群は交点二点の対を保存し、その補集合である
$Q$ 上の開軌道にも $L$ 上の開軌道にも推移的に作用し、二交点を交換する。
実際、二点を $0,\infty$ とすると対角群が両方の開軌道を与える。
接線の場合も一点の安定化群が $Q$ と接線それぞれの接点の補集合で推移的である
（標準二次 Veronese 表現で平行移動の作用により確認できる）。
三角形では頂点を座標点に取れば、対角群と座標の置換が表の二つの軌道を与える。
共点の三直線では、直線束の射影変換で任意の三方向を互いに移せる。
さらに、共点を $[0:0:1]$ とすると、
$[x:y:z]\mapsto[x:y:z+ax+by]$ が各直線上の共点の補集合で推移的である。
二直線の配置も、交点と二直線を座標位置に取れば、各直線上で交点の補集合が一軌道となる。
三重直線ではその安定化群が直線上で推移的である。これで各場合を尽くす。

以下の有理係数計算で非被約構造も含めた同型型を決定する。
$p=(0,0)$ とし、アフィン式を全次数3に斉次化して $C$ を復元する。
最後の行以外は、消去後の自由変数を必要に応じて $s,t,r$ と付け直した。
「環」欄には $\mathbb Q$ 上で計算した座標環を記す。$k=\mathbb Q$ とおく。
\[
\begin{array}{lll}
\text{型と点} & P(v,w)&\text{環}\\ \hline
\text{割線上のみ} &v(1+v^2+w^2)& k\\
\text{割線型、二次曲線上のみ} &v+w+vw& k\\
\text{割線型、交点} &v(w+v^2+w^2)& k[t]/(t^2-1)\\
\text{接線上のみ} &v(v+(w+1)^2)&k\\
\text{接線型、二次曲線上のみ} &(v-1)(v+w^2+2w)&k\\
\text{接線型、接点} &v(v+w^2)&k[t]/(t^2-1)\\
\text{三角形、辺} &v(1+w)(1+v)&k[t]\\
\text{三角形、頂点} &vw(1+v+w)&k[s,t]/(st)\\
\text{共点型、辺} &v(v+w-1)(w-1)&k[t]\\
\text{共点型、共点} &vw(v+w)&k[s,t]\\
2L+M,\ p\in L\setminus M & v^2(1+w)&k[s,t]/(t^2)\\
2L+M,\ p\in M\setminus L & v(1+w)^2&k[s,t]\\
2L+M,\ p\in L\cap M &v^2w&k[r,s,t]/(st)\\
3L &v^3& k[A,B,C',D]/(AD-2BC',B^2,BD)
\end{array}
\]
具体的には \eqref{p12:eq:cubic8} で $P$ の九係数を読み、
\eqref{p12:eq:normal8} の $H$ に代入する。このアフィン空間では
$J/F$ の十二係数が Jacobi イデアルを生成する。
まず $F$ 倍を除いても、すべての係数が生成するイデアルは変化しない。
実際 $F=w+uv$ は $w$ についてモニックなので、積の係数から商の係数を
最高 $w$ 次数から順に回収できる。六変数
$(x_0,x_1,y_0,y_1,z_0,z_1)$ の辞書式順序を用いる。
例えば割線型の交点の行のイデアルは
\[
 (x_0,\ x_1+z_1,\ 2y_0+z_1-1,\ y_1,\ z_0,\ z_1^2-1),
\]
三重直線の行では
\[
 (x_0z_1-2y_0z_0,\ x_1-2y_0,\ y_0^2,\ y_0z_1,\ y_1)
\]
である。後者では $A=x_0,B=y_0,C'=z_0,D=z_1$ とおけばよい。
全行について、元の十二係数と消去表示を与える生成系のそれぞれから被約
Gröbner 基底を $\mathbb Q$ 上で計算し、一致を照合する完全なコードを
\path{p1p2/verify_extension_fibers.py} に収録した。
$\mathbb C$ への基底変換で主張のスキームを得る。
\end{proof}


\subsection{del Pezzo 因子から得られる三つの既約成分}
以下の $dF$ などの外微分は各因子のアフィン錐上で計算する。
表示した1形式は多重次数が $-K_X=(2,3)$ で、両 Euler 場を消すので
大域切断 $\Omega^1_X(2,3)$ を与える。
また、積を外に出して書く logarithmic 1形式は閉形式なので可積分である\cite[\S3.3]{p12-Pym}。

\begin{theorem}\label{p12:th:threecomponents}
次の三つの族の Zariski 閉包は、$\operatorname{Pois}(X)$ の互いに異なる既約成分を与える。
\begin{enumerate}
\item $F\in H^0(\mathcal O(1,2))$、$U\in H^0(\mathcal O(1,0))$、
$L\in H^0(\mathcal O(0,1))$ に対し
\[
\alpha_{5}=UL\,dF-FL\,dU-2FU\,dL,
\qquad \dim C_{(1,2)(1,0)(0,1)}=14.
\]
一般の $\{F=0\}$ は5次 del Pezzo で、その上の Poisson 構造は $(UL)|_S$ に対応する。
\item $F\in H^0(\mathcal O(1,1))$、$U\in H^0(\mathcal O(1,0))$、
$Q\in H^0(\mathcal O(0,2))$ に対し
\[
\alpha_{8a}=2UQ\,dF-2FQ\,dU-FU\,dQ,
\qquad \dim C_{(1,1)(1,0)(0,2)}=12.
\]
一般の $\{F=0\}$ は8次 del Pezzo で、その上の Poisson 構造は $(2UQ)|_S$ に対応する。
\item $F,G\in H^0(\mathcal O(1,1))$、$L\in H^0(\mathcal O(0,1))$ に対し
\[
\alpha_{8b}=L(G\,dF-F\,dG),
\qquad \dim C_{(1,1)^2(0,1)}=11.
\]
一般の $\{F=0\}$ は8次 del Pezzo で、その上の Poisson 構造は $(LG)|_S$ に対応する。
\end{enumerate}
\end{theorem}
\begin{proof}
各表示は、非零な積切断で割るとそれぞれ
\[
\frac{dF}{F}-\frac{dU}{U}-2\frac{dL}{L},\qquad
2\frac{dF}{F}-2\frac{dU}{U}-\frac{dQ}{Q},\qquad
\frac{dF}{F}-\frac{dG}{G}
\]
となり、閉1形式である。従って $\alpha\wedge d\alpha=0$、すなわち $[\pi,\pi]=0$。
第一の族では $\alpha|_{F=0}=UL\,dF$、第二では $\alpha|_{F=0}=2UQ\,dF$、
第三では $\alpha|_{F=0}=LG\,dF$ であり、指定された曲面は Poisson 因子である。

成分性は次の線形代数の証明書で確立する。$\operatorname{Pois}(X)$ は34変数、30二次式で定義されるので、
ある Poisson 点で Jacobi の微分の階数が $34-d$ 以上、族をパラメーター表示する写像の微分の階数が $d$ 以上ならば、
族の閉包の局所次元は $d$ 以上、Poisson スキームの接空間次元は $d$ 以下となる。
よってその点は滑らかで、閉包は $d$ 次元の既約成分である。

アフィンチャート $([u:1],[v:w:1])$ における次の点を選ぶ。
\[
\begin{array}{c|l}
\mathrm{(5)}&F=uvw+(v-1)(w-1),\quad U=2u+3,\ L=v+2w+1\\
\mathrm{(8a)}&F=uv+w,\quad U=2u+3,\ Q=2v^2+vw+3w^2+5v+7w+11\\
\mathrm{(8b)}&F=(u+1)(v+2w+1)+(2u+3)(2v-w+3),\\
&G=(3u+1)(v+3w+2)+(u+4)(2v+w+5),\quad L=2v+3w+5.
\end{array}
\]
\eqref{p12:eq:jacobi} を用いて $30\times34$ の Jacobian 行列を、
全ての定義因子の係数で微分して、対応する parameter differential 行列を作る。
両行列は整数係数である。$\mathbb F_{101}$ へ還元した後、
それぞれ次の階数以上を持つことが、付属の独立した消去コードで確かめられる。
\[
\begin{array}{c|r|r|r}
\text{点}&\mathrm{rank}(dJ)&\mathrm{rank}(d\Phi)&34-\mathrm{rank}(dJ)\\\hline
\mathrm{(5)}&\ge20&\ge14&\le14\\
\mathrm{(8a)}&\ge22&\ge12&\le12\\
\mathrm{(8b)}&\ge23&\ge11&\le11
\end{array}
\]
全ての不等号は $\mathbb F_{101}$ 上で抽出した非零小行列式を有理数上に持ち上げることで得られる。
具体的には、行・列を辞書順に走査して独立な極大 minor を選ぶと、
Jacobi 行列の minor の行列式は mod $101$ で順に $71,20,54$、
パラメーター微分行列では $76,75,70$ となる。
各行列の作り方、列順、上記数値および必要な rank は同梱の
\path{p1p2/verify_del_pezzo.py} に完全に再現される。
この有限体上の非零 minor は特性0での非零を保証する。
よって $d=14,12,11$ として上記の成分性が得られる。

相異なる次元の成分は当然異なる。また11次元の $C_{(1,1)^2(0,1)}$ は
一般点で $\{L=0\}$ に沿って双ベクトルが零になる一方、
$C_{\mathrm{tor}}$ の一般点では零集合に余次元1の成分がない。
従って、この二つの11次元既約成分も互いに異なる。
\end{proof}

\begin{remark}
固定した $(1,1)$ と $(1,2)$ 因子への延長可能集合は定理\ref{p12:th:fixed8} と\ref{p12:th:extension5} で集合論的に決定した。
\end{remark}


\subsection{直線束と直線・二次曲線からの成分}

\subsubsection{射影平面の直線束からの成分}
$\Omega^1_X(2,3)=
(\Omega^1_{\PP^1}(2)\boxtimes\mathcal O_{\PP^2}(3))
\oplus(\mathcal O_{\PP^1}(2)\boxtimes\Omega^1_{\PP^2}(3))$
と分解する。
$H\in H^0(X,\mathcal O(2,1))$ と
$\omega\in H^0(\PP^2,\Omega^1_{\PP^2}(2))$ に対して
$\alpha=H\,\mathrm{pr}_2^*\omega$ は $\Omega^1_X(2,3)$ の切断である。
$\omega$ は平面から来るので $\omega\wedge d\omega=0$、従って
$\alpha\wedge d\alpha=0$ である。
\begin{theorem}\label{p12:th:fibration}
$\alpha=H\omega$ の全体の Zariski 閉包
$C_{\mathrm{fib}}$ は $\operatorname{Pois}(X)$ の11次元既約成分である。
その一般の切断は $(2,1)$ 型の零因子を持つ。
\end{theorem}
\begin{proof}
$H^0(X,\mathcal O(2,1))$ は9次元、
$H^0(\PP^2,\Omega^1(2))$ は3次元で、
パラメータ空間は既約である。
$\omega=dw$ および
\[
H=(1+u+u^2)+(2+3u+5u^2)v+(7+11u+13u^2)w
\]
に対応する点 $\pi=(H,0,0)$ では、
Jacobi 方程式の微分の階数が少なくとも23、
族のパラメータ微分の階数が少なくとも11である。
実際、34変数を本文の基底順に並べ、30式の微分から
$23\times23$ の非零小行列式を、族の微分から $11\times11$ の
非零小行列式を抽出すると、それぞれ
$\mathbb F_{101}$ 上で行列式 $6,94$ が得られる。
再現スクリプト \path{p1p2/verify_new_components.py} により
有限回の整数演算として検証できる。
従ってこの点は Poisson 概型の滑らかな11次元点であり、
族の閉包は既約成分である。
一般の $H$ は $(2,1)$ 型の既約因子を定め、$\omega$ の零集合は
$\PP^2$ の一点だから、共通の因子として $H$ が現れる。
\end{proof}

\subsubsection{直線と滑らかな二次曲線からの成分}
$L\in H^0(\PP^2,\mathcal O(1))$ と
$Q\in H^0(\PP^2,\mathcal O(2))$ を取り、
$h\in H^0(\PP^1,\mathcal O(2))$ および $a\in\mathbb C$ に対して
\begin{equation}\label{p12:eq:lineconic}
\alpha=aLQ\,du+h(u)(L\,dQ-2Q\,dL)
\end{equation}
を定める。$du$ はアフィン座標で表した
$\Omega^1_{\PP^1}(2)\simeq\mathcal O_{\PP^1}$ の標準生成切断である。
分子 $L\,dQ-2Q\,dL$ は Euler 縮約が零の
$\Omega^1_{\PP^2}(3)$ の切断なので、
\eqref{p12:eq:lineconic} は大域切断を定める。
$LQ\ne0$ では
\[
\frac\alpha{LQ}=a\,du+h(u)\,d\log\!\left(\frac Q{L^2}\right),
\]
であり、これは Frobenius 条件を満たす。
\begin{theorem}\label{p12:th:lineconic}
\eqref{p12:eq:lineconic} の全体の Zariski 閉包
$C_{\mathrm{line\text{-}conic}}$ は $\operatorname{Pois}(X)$ の
11次元既約成分である。一般の切断には余次元1の零成分がなく、
$C(v,w)$ は滑らかな二次曲線と直線の積を表す。
\end{theorem}
\begin{proof}
$L$、$Q$、$h$、$a$ の選択空間は既約である。
アフィン座標において
\begin{align*}
 L&=2v+3w+5,\\
 Q&=2v^2+vw+3w^2+5v+7w+11,\\
 h&=1+2u+3u^2,\qquad a=7
\end{align*}
と取る。対応する $\pi$ の Jacobi 方程式の微分階数は
$\mathbb F_{101}$ 上で23、
13個のパラメータに関する微分階数は11である。
抽出した極大小行列式はそれぞれ $6,22\pmod{101}$ で非零であり、
有理数体上でも同じ階数下界が成り立つ。
定理\ref{p12:th:fibration} と同様に11次元既約成分となる。

$C=\alpha_u=aLQ$ は直線と滑らかな二次曲線の積である。
共通の因子が存在すれば $L$ または $Q$ が
$\alpha_v,\alpha_w$ をも割らねばならないが、
$L=0$ では $-2hQ\,dL$、$Q=0$ では $hL\,dQ$ となり、
一般の $L,Q$ ではどちらも恒等的には零でない。
従って一般の $\alpha$ は余次元1の零成分を持たない。
\end{proof}

\subsection{四因子の対数形式と成分の交わり}
\label{p12:sec:four}
$F,U,L,M$ の双次数を順に $(1,1),(1,0),(0,1),(0,1)$ とする。
$a,b\in\mathbb C$ に対し
\begin{equation}\label{p12:eq:four}
 \alpha=FULM\left(a\,d\log\frac{F}{UM}
                    +b\,d\log\frac L M\right)
\end{equation}
を考える。両商は次数零であり、この式は大域的な捻れ1形式に延びる。
非零の因子積で割ると閉形式になるので可積分である。
\begin{theorem}\label{p12:th:four}
\eqref{p12:eq:four} の像の閉包 $C_4$ は $\operatorname{Pois}(X)$ の12次元既約成分である。
$a\ne0$ のとき $F=0$ は Poisson 因子で、そこでの制限は
$a(ULM)|_{F=0}$ に対応する。$F$ が滑らかならこの因子は8次 del Pezzo 曲面である。
\end{theorem}
\begin{proof}
パラメータ空間は既約である。一般の四因子の配置では
$U=1,\ L=v,\ M=w,\ F=u+v+w$ をアフィン表示とできる。
実際、$U$ を無限遠に置き $F=uD+E$ とする。平面の三つの一次形式 $L,M,D$
を座標に取り、$D=1$ 上で $u$ の平行移動と $v,w$ の定数倍を用いる。
この配置で $a=1,\ b=r$ と取ると、係数は
\[
 A=vw-(1+r)vF,\qquad B=-vw-rwF,\qquad C=vw.
\]
$r=2$ における30行34列の Jacobi 微分の階数は22であり、
族の微分の階数は12である。いずれも有理数上の正確な消去である。
後者の列は、各 $F,U,L,M$ をその切断空間の単項式基底の元で
一つずつ置換したものと $(a,b)$ の二つの基底変分で尽くされる。
この計算と下記の不変因子は \path{p1p2/verify_residues.py} で再現できる。
従って像の閉包の次元と接空間次元はこの点で一致する。
$F=0$ 上では $\alpha=aULM\,dF$ であり、制限の主張も従う。
\end{proof}

12次元の既出成分との相異性は $C$ の因数分解でわかる。
\eqref{p12:eq:four} の $du$ 係数はスカラーを除き
$LM(U\,\partial_u F-F\,\partial_u U)$ で、つねに完全分解する。
定理\ref{p12:th:threecomponents} の第2族では一般に滑らかな二次曲線が因子となる。
したがってこの二つの12次元成分は異なる。

留数比に関する跳躍も決定できる。上記一般配置を固定し、
$r=b/a$ と置くと、$dJ$ の $\mathbb Q[r]$ 上の非零 Smith 不変因子は単元倍を除いて
\[
 \underbrace{1,\ldots,1}_{19},\quad f,f,\quad
 f\,r(r+1)(2r+1),\qquad f=(r-1)(r+2).
\]
$r=\infty$ では $(a,b)=(0,1)$ を代入する。これにより以下を得る。
右列には各値で実際に交わる成分を一つ記した。
\[
\begin{array}{c|c|l}
 r &\dim T_\pi\operatorname{Pois}(X)& \text{所属を与える表示}\\\hline
 1,-2&15&C_{(1,2)(1,0)(0,1)}\quad (F L)/(U M^2)\ \text{または }(F M)/(U L^2)\\
 -1/2&13&C_{(1,1)(1,0)(0,2)}\quad F^2/(U^2 L M)\\
 0,-1&13&C_{(1,1)^2(0,1)}\quad F/(U M)\ \text{または }F/(U L)\\
 \infty&14&C_{\mathrm{fib}}\quad U F (M\,dL-L\,dM)\\
 \text{上記以外}&12& C_4\text{ の滑らかな点}
\end{array}
\]
例えば $r=1$ では定理\ref{p12:th:threecomponents} 第1族で
$(1,2)$ 型の分子を $FL$ に特殊化した形になる。
最後の行では12次元成分上にあることと接空間次元から滑らかさが従う。

同じ一般配置で $a\ne0$ とし、$x=L/M$, $y=F/(UM)$, $z=w$ と置くと、
$\mathbb C(X)=\mathbb C(x,y,z)$ で、葉層を $d\log y+r\,d\log x$ が定める。
$r\notin\mathbb Q$ なら有理第一積分は存在しない。実際、第一積分は
$\partial_z$ で消えるから $\mathbb C(x,y)$ に属し、
$x\partial_x-r y\partial_y$ でも消える。互いに素な分子・分母は
それぞれこの導分の固有多項式で、固有値は等しい
（整除と全次数から従う）。異なる単項式は異なる重みを持つので、
比は定数である。


\subsubsection{成分の相異性と混合階数1}
\begin{proposition}\label{p12:prop:distinct}
ここまでに構成した $C_{\mathrm{Ric}}$ 以外の八つの既約成分は互いに異なる。
\end{proposition}
\begin{proof}
次元が異なれば成分は異なる。12次元の二成分は第\ref{p12:sec:four}節で区別した。11次元の四成分は次の幾何的不変量で区別する。
$C_{\mathrm{fib}}$ のすべての点では
$C(v,w)\equiv0$ である。他の三つの11次元成分では一般点で $C\ne0$。
$C_{\mathrm{tor}}$ の一般点では $C$ は相異なる三直線の積で、
この性質の閉包は完全分解可能な三次式の錐の中にある。
一方 $C_{\mathrm{line\text{-}conic}}$ の一般点の $C$ は直線と
既約で滑らかな二次曲線の積であるから、両者は一致しない。
$C_{(1,1)^2(0,1)}$ は $\alpha=L\alpha'$ と書けるため、
$\mathcal O(0,1)$ 型の余次元1零因子をもつ切断の閉集合に含まれる。
対して $C_{\mathrm{line\text{-}conic}}$ の一般点は
定理\ref{p12:th:lineconic} の通り余次元1の零因子を持たない。
従って後者と pencil 成分も異なる。
一般のトーラス双ベクトルの零点集合は余次元2以上である。実際、開トーラス軌道と各境界因子の開軌道では接方向が少なくとも二つ残り、一般の係数なら外積は零にならない。従ってトーラス成分と pencil 成分も異なる。
\end{proof}

\begin{proposition}\label{p12:prop:rankone}
$X$ 上の大域双ベクトル
\[
 \pi=\operatorname{pr}_2^*\beta_C+h(u)\,\partial_u\wedge V,
 \qquad \beta_C=C(v,w)\partial_v\wedge\partial_w,
\]
を考える。$h(u)\partial_u$ は $\PP^1$ 上の非零大域ベクトル場、
$V=P(v,w)\partial_v+Q(v,w)\partial_w$ は $\PP^2$ 上の大域ベクトル場とする。
このとき $\pi$ が Poisson である必要十分条件は
\[
\mathcal L_V\beta_C=0,
\quad\text{すなわち}\quad
 P\,C_v+Q\,C_w-C(P_v+Q_w)=0
\]
である。
\end{proposition}
\begin{proof}
\eqref{p12:eq:jacobi} に $A=hP,\ B=hQ$ を代入すると、
$h'$ に関する二項が打ち消され、
$J=h(P C_v+Q C_w-C(P_v+Q_w))$ となる。
$h\ne0$ ならこの恒等式が示す条件に同値である。
\end{proof}

\subsection{直線と滑らかな二次曲線を射影三次曲線とする場合}
\label{p12:sec:conic-stratum}
$Y=\PP^2$ とし、その Poisson 双ベクトル $G$ の零因子が直線 $L$ と滑らかな
二次曲線 $Q$ の和である場合を考える。$L$ が $Q$ に接する場合も許す。
平面に沿う射影を $G$ に固定する。大域ベクトル場 $V_i$ により
\[
 \pi=G+V(t)\wedge\partial_t,\qquad
 V(t)=V_0+tV_1+t^2V_2
\]
と表す。この符号では \eqref{p12:eq:jacobi} は
$\mathcal L_{V(t)}G=V(t)\wedge V'(t)$ である。$V_i$ がすべて $Q$ に
接することを「$Q$ に沿って対数的」と呼ぶ。これは対応する平坦接続の
すべての成分が $Q$ の一般点で対数極を持つという条件（正則なものも含む）
に一致する。$L$ に沿う対数性も同様に定める。

\begin{proposition}\label{p12:prop:one-bad-conic}
すべての $\pi$ は $L,Q$ の少なくとも一方に沿って対数的である。
\end{proposition}
\begin{proof}
$Y\setminus(L\cup Q)$ の基本群は二次曲線の正のメリディアン $m_Q$ で生成される
無限巡回群で、直線のメリディアンは $m_L=m_Q^{-2}$ である。
実際、$L$ を無限遠とすると、接線の場合はアフィン二次曲線を $y=x^2$ と書け、
補集合は $\mathbb C\times\mathbb C^*$ である。割線の場合は二次曲線を $xy=1$
と書ける。補集合からさらに $x=0$ を除くと座標 $(x,xy)$ により
$\mathbb C^*\times(\mathbb C\setminus\{1\})$ を得る。
この開集合は元の補集合の基本群へ全射し、第1因子のループは $y=0$ 上で縮む。
残る生成元の位数は $xy-1$ の巻き数により無限である。いずれの場合も
有理関数 $Q/L^2$ の巻き数がメリディアンの関係を与える。

$Q$ に沿って非対数的なら、対応するトレース零の階数2平坦接続で
$m_Q$ のモノドロミー固有値はともに $-1$ である。実際、一般点の
局所座標 $z,s$（$Q=(z=0)$）で接続行列を $(A\,dz+B\,ds)/z$ と書くと、
平坦性から $[A|_{z=0},B|_{z=0}]=B|_{z=0}\ne0$ を得る。
従って $A|_{z=0}$ の固有値は $\pm1/2$ であり、横断円板上の正則特異接続の
モノドロミーの固有値はともに $-1$ である。
このとき $m_L=m_Q^{-2}$ の固有値はともに $+1$ であるから、
同じ理由で $L$ に沿う非対数極は存在しない。
\end{proof}

\begin{theorem}\label{p12:th:conic-log-classification}
$G$ を固定し、$Q$ に沿って対数的な $\pi$ を考える。混合テンソルの階数が
$1$ 以下の場合は $C_{\mathrm{line\text{-}conic}}$ の表示
\eqref{p12:eq:lineconic} から得られる。階数が $2$ 以上の場合はすべて
定理\ref{p12:th:threecomponents} の第2族 $C_{(1,1)(1,0)(0,2)}$ の表示から得られ、
階数はちょうど $2$ である。
\end{theorem}
\begin{proof}
$W=\operatorname{Lie}\operatorname{Aut}(Y,Q)\cong\mathfrak{sl}_2$ とする。
$W$ の外積は $Q$ 上で零であるから
\[
 j:\textstyle\bigwedge^2 W\longrightarrow Q\cdot H^0(Y,\mathcal O(1))
\]
を定める。$\operatorname{Aut}(Y,Q)$ は $Q$ 外で開軌道を持つので
$j$ は零写像でない。両辺とも $\operatorname{PGL}_2$ の3次元既約表現である。
実際 $Y$ を二変数二次形式の射影平面と同一視すれば $Q$ は判別式二次曲線であり、
$H^0(Y,\mathcal O(1))$ は二次対称冪である。従って $j$ は同型である。
$k=[\ ,\ ]\circ j^{-1}$ とおき $h=k(G)\ne0$ とする。

仮定は $V_i\in W$ である。$k$ の同変性と $k(V\wedge V')=[V,V']$ から
Jacobi 方程式は
\[
 [V_0,h]=[V_0,V_1],\quad [V_1,h]=2[V_0,V_2],\quad
 [V_2,h]=[V_1,V_2]
\]
となる。階数1以下なら $V_i$ はすべて $h$ と可換である。
$\mathfrak{sl}_2$ の非零元の中心化代数は一次元なので $V(t)=f(t)h$ であり、
これは \eqref{p12:eq:lineconic} に一致する。実際 $\mathcal L_hG=0$ で、
$L\,dQ-2Q\,dL$ に対応する非零ベクトル場は $W$ に属し $G$ を保つ
（$G$ で割ると $d\log(Q/L^2)$ に比例する）。

階数2以上なら、二つの独立な値を $0,\infty$ で取るように
$\PP^1$ の座標を選べるから、$V_0,V_2$ は独立としてよい。第1、第3式より
$h-V_1=aV_0,\ h+V_1=bV_2$ であり、第2式より $ab=-4$ である。
従って $V(t)=(1-at/2)V_0+t(t+b/2)V_2$ は $t=2/a$ で零となる。
この点を無限遠に移せば、新座標の二次係数は零で、第2式は
$[V_1,h]=0$ である。$V_0,V_1$ は独立なので第1式から $V_1=h$ を得る。
こうして標準形は $V(t)=Z+th$, $Z\in W\setminus\mathbb Ch$ となる。

最後にこの標準形を対数族と照合する。$Q$ の方程式を固定し、
$G$ の式が $2QL$ となるように直線の方程式 $L$ を選ぶ。
定理\ref{p12:th:threecomponents} 第2族で $U$ を無限遠、
$F=tL+M$（$M$ は任意の平面一次形式）とすると、アフィン表示は
\[
 2QL\,dt+(2Q\,dM-M\,dQ)+t(2Q\,dL-L\,dQ)
\]
となる。$M\mapsto 2Q\,dM-M\,dQ$ に対応するベクトル場は $Q$ に接する。
この線形写像は単射である。非零の $M$ が核にあれば
$d(M^2/Q)=0$ となり、$Q$ が一次形式の平方となるためである。
従ってこの写像は $W$ 上への同型である。$M$ を $L$ と独立に選べば階数は2で、
既に示した方程式から一次係数は $h$ である。この係数は $M$ と独立であり、
定数係数は $W$ の任意の元を取るので、すべての標準形が得られる。
\end{proof}

\begin{theorem}\label{p12:th:secant-classification}
$L$ が $Q$ と横断的に交わるとする。射影双ベクトル $G$ を固定した
集合論的ファイバーは、\eqref{p12:eq:lineconic} および
定理\ref{p12:th:threecomponents} の三つの表示で $G$ を誘導するものの和集合である。
特に、三次曲線が一般の「滑らかな二次曲線とその割線」であるような
Poisson 既約成分は、これらの四成分に尽くされる。
\end{theorem}
\begin{proof}
命題\ref{p12:prop:one-bad-conic}と定理\ref{p12:th:conic-log-classification}により、
$L$ に沿って対数的で $Q$ に沿って非対数的な場合を調べればよい。
$L$ を無限遠、$Q=(q=xy-1=0)$ ととる。$L$ に接する射影ベクトル場は
$x,y$ に関しアフィン一次式を係数にもつ。従って $G$ を係数とする接続行列は
\[
 \Omega=\frac{(a+bx+cy)\,dx+(d+ex+fy)\,dy}{q},
 \qquad a,b,c,d,e,f\in\mathfrak{sl}_2.
\]
ここで分子の定数倍に $G$ の尺度を吸収した。平坦性
$d\Omega+\Omega\wedge\Omega=0$ の六つの係数は
\begin{align*}
 [e,b]=b,\qquad [c,f]=f,\qquad [b,f]+[c,e]=0,\\
 a+[a,e]+[b,d]=0,\qquad -d+[a,f]+[c,d]=0,\qquad [a,d]=e-c .
\end{align*}
$Q$ の一般点で $x$ を一定にした横断円板上の留数は
$e+d/x+f/x^2$ である。前命題の局所計算からその固有値は
$\pm1/2$ である。行列式の恒等式で $x\to\infty$ とすると
$e$ も同じ固有値を持つ。$x,y$ を交換すれば $c$ についても同様である。

$H=\operatorname{diag}(1/2,-1/2), J=E_{12}, K=E_{21}$ とする。
$b,f$ の非零なものはそれぞれ $\operatorname{ad}e,\operatorname{ad}c$ の
重み1のベクトルなので冪零である。両方が非零で比例しないなら、
それぞれの核の直線を基底として $b=BJ,f=DK$ と書ける。
最初の二式から $e=H+rJ,c=-H+sK$ であるが、第3式の
$J,K,H$ 成分から順に $r=s=0,BD=0$ となり矛盾する。
両方が非零で比例する場合は $b=BJ,f=DJ$ とできる。このとき
$e=H+rJ,c=H+sJ$ で、第3式から $r=s$ である。
定数共役により $c=e=H$ としてよい。

一方だけ非零なら $x,y$ の交換で $b=BJ\ne0,f=0$ とできる。
$e=H$ としたとき第3式と $c$ の固有値から $c=\pm H$ となる。
$c=-H$ なら第5式より $d=d_-K$、第6式より $d_-\ne0$ かつ
$a$ の $H$ 成分は零である。第4式の $H$ 成分 $2Bd_-=0$ はこれと矛盾する。
$b=f=0$ の場合も $[c,e]=0$ より $c=\pm H$, $e=H$ とできる。
このうち $c=-H$ のとき、第4--6式から
$a=AJ$, $d=A^{-1}K$ ($A\ne0$) を得る。

残った $c=e=H$ の場合は $b=BJ,f=DJ$（零も可）であり、第4式の
$K$ 成分から $a_K=0$、第5式の $K$ 成分から $d_K=0$ を得る。
両式の $H$ 成分から $a_H=d_H=0$ となるので $a,d\in\mathbb CJ$ である。
以上により定数共役の後に接続行列は次の二形に限られる。
\begin{align*}
 &\Omega=(H\,dq+J\,dP)/q,
    &&P=p_1x+p_2y+p_3 x^2+p_4 y^2,\\
 &\Omega=\big[(AJ-Hy)\,dx+(A^{-1}K+Hx)\,dy\big]/q,
    &&A\ne0.
\end{align*}
第1形の射影水平方程式は $d(qt+P)=0$ である。
$F=qt+P$ を斉次化すると $(1,2)$ 型となり、得られる双ベクトルは
定理\ref{p12:th:threecomponents} の第1族（$U,L$ が無限遠）に属する。
第2形は
\[
 M=\begin{pmatrix}y&-A\\-A^{-1}&x\end{pmatrix},\qquad
 \det M=q,\qquad
 \Omega=M^{-1}dM-\tfrac12\operatorname{tr}(M^{-1}dM)I
\]
である。従って $M(t,1)^{\mathsf T}$ の成分比が第一積分である。
この二成分の斉次化はともに $(1,1)$ 型であり、
分母を払った1形式に無限遠の $L$ を掛けたものが
定理\ref{p12:th:threecomponents} の第3族である。
実際、無限遠を除くチャートでは両辺の $dt$ 係数は
$q$ の定数倍で等しく、正規化した水平方程式も同じであるから、
$G$ の尺度を合わせれば双ベクトルとして一致する。
\end{proof}

\subsubsection{直線が二次曲線に接するときの非対数的な極}
\label{p12:sec:tangentgoodline}
\begin{proposition}\label{p12:th:tangent-line}
平面三次曲線 $L+Q$ が滑らかな二次曲線 $Q$ とその接線 $L$ からなり、
平坦接続が $L$ の一般点で対数的、$Q$ の一般点で非対数的であるとする。
対応する Poisson 構造は $C_{(1,2)(1,0)(0,1)}$ 型の族または
$C_{(1,1)^2(0,1)}$ 型の族に属する。
\end{proposition}
\begin{proof}
平面の射影変換と全体のスカラー倍を施し、
$L$ を無限遠直線、$Q$ のアフィン方程式を $q=y-x^2$ とする。
$L$ で対数的であることは混合項の各射影ベクトル場が $L$ に接することと同値であり、
それらはアフィン線形な係数を持つ。従って水平持ち上げを
$\partial_i+\Omega_i$ と書く規約では、
\[
 \Omega=(X\,dx+Y\,dy)/q,\qquad
 X=Mx+Ny+A,\quad Y=Px+Ry+B,
\]
ここで係数は $\mathfrak{sl}_2$ に属する。
$\mathfrak{sl}_2$ を射影直線のベクトル場として用いるときもその Lie 括弧を保つ。
平坦性を
$q(\partial_x Y-\partial_y X)-q_xY+q_yX+[X,Y]=0$
と展開すると、次の六式を得る。
\begin{align*}
 [M,P]+P+N&=0, & [M,R]+[N,P]+2R&=0, & [N,R]&=0,\\
 (1-\operatorname{ad}B)M+2B+[A,P]&=0, &
 P+[N,B]+[A,R]&=0, & [B,A]&=A.
\end{align*}
$(x,q)$ 座標では $dq/q$ の係数を $Q$ に制限した行列は
$Y(x,x^2)=B+xP+x^2R$ である。$dx/q$ の係数の $Q$ 上の制限が
一般点で非零なので、平坦性の二位の極の係数から
$\operatorname{ad}Y(x,x^2)$ は固有値1を持つ。従って全ての $x$ に対し
$\operatorname{tr}(Y(x,x^2)^2)=1/2$ である。
特に定数共役で $B=h=\operatorname{diag}(1/2,-1/2)$ とできる。
上・下三角の標準行列をそれぞれ $e,f$ とし、
$[h,e]=e,\ [h,f]=-f,\ [e,f]=2h$ とする。
先の六式の最後から $A=ae$、トレース恒等式の一次係数から
$P=be+cf$ と書ける。$R=E e+K h+D f$ とおくと、その二・三・四次係数は
\[
 K=-2bc,\qquad bD+cE=0,\qquad K^2+4ED=0
\]
となる。六式の下段の最初の二式、ついで上段の第一式を比較すると
\begin{align*}
 M&=m e+H h, & H&=-2(1+ac),\\
 N&=(b-aK)e+n h-c f, & aD&=0,\qquad Hc=0,\qquad n=-2mc.
\end{align*}

$c=0$ なら $K=0,\ H=-2,\ N=P=be$ である。
$2R+[m e-2h,R]=0$ の $f$ 成分から $D=0$、従って $R=Ee$ となる。このとき
\[
 \Omega=h\,dq/q+e\,dg/q,\qquad
 g=m x^2/2+bxy+ax+E y^2/2 .
\]
基底 $h,e,f$ を $t\partial_t,t^2\partial_t,-\partial_t$ で実現すると、
水平な方向では $d(q/t+g)=0$ となる。$q$ と $g$ を
平面上の二次式に斉次化すれば、この有理第一積分は
$F/(U L^2)$ の形をとる。ここでアフィン座標では $F=q+tg$、$U=t$ であり、
その双次数は $(1,2),(1,0)$ である。
実際 $U L\,dF-FL\,dU-2FU\,dL$ のアフィン式は
$-q\,dt+t\,dq+t^2\,dg$ で、水平な1形式の定数倍に一致する。
従って主張の第一の族に入る。

$c\ne0$ なら $H=0,\ ac=-1$ である。$aD=0$ から $D=0$、
トレース条件から順に $K=b=E=0$、従って $R=0$ となる。
$[N,P]=0$ から $m=0$ が従い、$M=0$, $N=-P=-cf$ を得る。
対角共役で $a=1,c=-1$ とする。同じベクトル場の実現では
水平な1形式は
\[
 q\,dt+(y-t^2)\,dx-(t+x)\,dy
 = G\,dF-F\,dG,\qquad F=t+x,\quad G=tx+y
\]
である。$F,G$ を双次数 $(1,1)$ に斉次化し $G\,dF-F\,dG$ に
$L$ を掛けることで元の大域双ベクトルを得る。
\end{proof}


\begin{theorem}\label{p12:th:wholeconic}
$C=LQ$ が直線と滑らかな二次曲線の積なら、$q^{-1}(C)$ の閉点を尽くすのは
$C_{\mathrm{line\text{-}conic}}$、$C_{(1,1)(1,0)(0,2)}$、
$C_{(1,2)(1,0)(0,1)}$、$C_{(1,1)^2(0,1)}$ の四族である。
\end{theorem}
\begin{proof}
命題\ref{p12:prop:one-bad-conic} による二分のうち、
二次曲線に沿って対数的な場合は定理\ref{p12:th:conic-log-classification} に従う。
残りは割線の場合は定理\ref{p12:th:secant-classification}、接線の場合は命題\ref{p12:th:tangent-line} に従う。
\end{proof}

\subsection{三本の非共点直線の上のファイバー}\label{p12:sec:triangle}
\begin{theorem}\label{p12:th:triangle}
射影三次曲線が三本の相異なる非共点直線の和である Poisson 双ベクトルは、
$C_{\mathrm{tor}}$, $C_4$, $C_{(1,2)(1,0)(0,1)}$ のいずれかに入る。
より正確には、三直線のうち平坦接続が一般点で非対数的なものの本数は
$0,1,2$ のいずれかであり、各場合に順に上記の成分に属する。
\end{theorem}
\begin{proof}
補集合は $(\mathbb C^*)^2$ で、三本のメリディアンは可換で積が1である。
非対数極のメリディアンのモノドロミー固有値はともに $-1$ である。
可換な行列は同時三角化できるので、この固有値条件が三本すべてでは成立しない。

まず少なくとも二本に沿って対数的とし、それらを $x=0,y=0$、残りを無限遠とする。
両軸に接する射影ベクトル場は
$x(a+\ell)\partial_x+y(b+\ell)\partial_y$ と書ける。
ここで $a,b$ は定数、$\ell$ は斉次一次式である（射影ベクトル場の
二次部分は共通の一次式と Euler 場の積である）。
$G$ のアフィン方程式は $xy$ の定数倍なので、接続行列は定数行列
$A,B,S,T\in\mathfrak{sl}_2$ を用いて
\[
 \Omega=(A+Sx+Ty)\,d\log x+(B-Sx-Ty)\,d\log y
\]
となる。この表示において無限遠でも対数的となる条件は $S=T=0$ である。
実際 $d\log(x/y)$ は無限遠直線に接する一般の方向への制限が非零で、
$Sx+Ty$ がその一般点に極を持たないためには $S=T=0$ が必要である。
平坦性は係数比較により
\[
 [A,B]=0,\qquad [S,A+B]=S,\qquad [T,A+B]=T
\]
と同値である。従って $S=T=0$ なら $A,B$ は共通の
$\mathfrak{sl}_2$ 元のスカラー倍である。この元が半単純の場合は
トーラス上の対数双ベクトルであり、冪零の場合も共通の元を半単純元から
特殊化できるので $C_{\mathrm{tor}}$ に属する。

$(S,T)\ne(0,0)$ なら $A+B$ の随伴表現は固有値 $-1$ をもつから、
定数共役で $A+B=-H$, $H=\operatorname{diag}(1/2,-1/2)$ とできる。
$A,B$ は $A+B$ と可換であり $A=aH,B=bH$, $a+b=-1$ である。
また $S=sJ,T=rJ$ ($J=E_{12}$) である。$\eta=a\,d\log x+b\,d\log y$
と書くと射影水平方程式は
\[
 dt+t\eta+(sx+ry)(d\log x-d\log y)=0 .
\]
一般の $a,b$ には $a(1+a)\ne0$ が成立する。このとき
$P=p x+k y$ を $(1+a)p=s,\ ak=r$ で定めると
$dP+P\eta=(sx+ry)(d\log x-d\log y)$ であり、
方程式は $d(t+P)+(t+P)\eta=0$ に一致する。
これは $t+P$ の斉次化を $(1,1)$ 因子とする四因子対数形式
\eqref{p12:eq:four} である。$a=0,-1$ についても、
固定した $s,r$ のもとで上の行列表示の $a$ を特殊化すれば $C_4$ に属する。

次に非対数極が二本の場合、それらを座標軸 $x=0,y=0$ に取る。
無限遠で対数的なので、$q=xy$ として
\[
 \Omega=\big[(a+bx+cy)\,dx+(d+ex+fy)\,dy\big]/q
\]
と書く。定理\ref{p12:th:secant-classification} の $xy-1$ の計算と比較すると、
平坦性の六式のうち定数式だけが $[a,d]=0$ に変わる。
$y=0$ と $x=0$ を横切る円板上の留数はそれぞれ $e+d/x$, $c+a/y$ であり、
その固有値が $\pm1/2$ であることから、$e,c$ にも同じ固有値がある。
$H,J,K$ を先の証明と同じ行列とする。
$b,f$ が非比例の場合の排除と、$b= B J, f=D J, c=e=H$ への帰着
（両方非零かつ比例する場合）は最初の三式だけを用いたので、そのまま適用できる。
一方だけ非零なら対称性から $b=BJ\ne0,f=0$ とでき、$e=H,c=\pm H$ である。
$b=f=0$ でも $e=H,c=\pm H$ とできる。
$c=-H$, $b=BJ\ne0,f=0$ のとき $d=d_-K$,
$a_K=0,\ a_H=-2Bd_-$ である。
$[a,d]=0$ のもとで $d_-\ne0$ なら $a\in\mathbb CK$ となり
$a_H=0$ と矛盾する。従って $d=0$ となるが、
このとき $x=0$ 上で $dy$ の分子 $d+fy$ が零となり対数的である。
$c=-H$, $b=f=0$ のときは $a=AJ,d=DK$ となり $[a,d]=0$ から
$AD=0$ を得るので、やはりどちらかの軸で対数的になる。

残る $c=e=H$ では、先の計算の定数式以外から
$a,b,d,f\in\mathbb CJ$ を得る。
従って $\Omega=(H\,dq+J\,dP)/q$（$P$ は二次以下の多項式）であり、
$q t+P$ が有理第一積分である。$qt+P$ を斉次化した $(1,2)$ 因子と
無限遠の二直線を用いた定理\ref{p12:th:threecomponents} 第1族の表示を得る。
\end{proof}

\subsection{共点三直線の Riccati 成分}
中心 $p$ で交わる三本の相異なる直線の和を $D$ とする。
$p$ を $y$ 軸方向の無限遠点に置き、三直線の方程式を
$f(x)=\prod_{i=1}^3(x-r_i)=0$ と取る。平面双ベクトルのスカラー倍は
以下の族に含める。ファイバー座標 $t$ に関する二次以下の多項式
$h_0,h_1$ を選び、
\begin{equation}\label{p12:eq:concurrent-family}
 \pi= f(x)\,\partial_x\wedge\partial_y+
 h(t,x)\,\partial_t\wedge\partial_y,\qquad h=h_0(t)+x h_1(t)
\end{equation}
とする。これは式~\eqref{p12:eq:jacobi} を満たす。この族を座標変換で動かした閉包を
$C_{\mathrm{Ric}}$ とする。

\begin{theorem}\label{p12:th:concurrent-riccati}
$C_{\mathrm{Ric}}$ は12次元の既約成分である。その一般閉点で
$\operatorname{Pois}(X)$ の接空間は13次元であり、この成分において $\operatorname{Pois}(X)$ は
生成点で非被約である。
\end{theorem}
\begin{proof}
共点三直線を与える非零の平面双ベクトルの層は、中心 $p$ の2次元と、
$p$ を消す二つの一次式についての binary cubic の4次元からなる。
中心は一意である。$p$ を固定し、同じ中心の直線束の各葉に接する
ベクトル場は $W_p$（命題~\ref{p12:prop:zero-cubic}）をなし、
今の座標では $(a+bx+cy)\partial_y$ と書ける。このような混合項では
Jacobi 式は $-f\partial_y(a+bx+cy)=0$（$a,b,c$ は $t$ の二次以下の式）なので
$c=0$ である。従って~\eqref{p12:eq:concurrent-family} は、この層上で
各 $V_i$ が三直線に接する族の全体であり、ファイバーは6次元線形空間である。
実際、三直線を保存する射影変換の単位成分は、交点 $p$ を固定し、
直線束上で三点を固定するので束に自明に作用する。
したがって族は既約な12次元である。

この族が既約成分となることをモノドロミーで示す。
式~\eqref{p12:eq:concurrent-family} の接続は、直線束の底上の
\[
 d+\lambda(M_0+x M_1)\,dx/f(x)
\]
の引き戻しである。ここで $M_0,M_1\in\mathfrak{sl}_2$ は各 $h_i\partial_t$
に対応し、$\lambda$ は混合項を動かすスカラーである。
$[M_0,M_1]\ne0$ とする。$r_i$ の周りの基点付きメリディアンのホロノミーは
$I-2\pi i\lambda R_i+O(\lambda^2)$、$R_i=(M_0+r_iM_1)/f'(r_i)$ である。
異なる二つの $R_i$ の括弧は非零なので、十分小さな $\lambda\ne0$ で
対応するホロノミー行列は可換でない。一方、三つのメリディアンのトレースは
いずれも $2$ に近く、$-2$ とは異なる。

上の一つの点の近傍で、三次曲線の被約性は保たれる。
有限個の基点付き閉路を固定すると、その像を避ける近傍でホロノミーが
パラメータに連続に依存するので、二つの行列の非可換性も保たれる。
被約平面三次曲線のうち、補集合の基本群が非可換なのは共点三直線だけである
\cite[\S2.3.2, (3a)--(3e)]{p12-Cogolludo}。
よってこの近傍のすべての Poisson 構造の平面成分は共点三直線である。
さらに近傍を縮めれば三本を局所的に番号付けでき、上で選んだ各メリディアンは
動いた直線の対応するメリディアンに自由ホモトピックである。
実際、各直線の滑らかな点に置いた小さい横断円板を使えばよい。
定理~\ref{p12:th:irreducible-cubic} の曲率計算により、ある直線に沿って
接続が非対数的なら、その一般のメリディアンのホロノミーのトレースは $-2$ である。
この局所的結論には基本群に関する仮定を必要としない。
トレースの連続性により三本とも対数的でなければならない。
従って $V_i$ はすべてこの三直線に接し、近傍の閉点は先に記述した
12次元族に属する。これより $C_{\mathrm{Ric}}$ は既約成分である。

次に接空間を調べる。固定 $f$ に対する六つの混合方向があるほか、
射影 $q$ の微分の像は $x$ の三次以下の多項式全体を含む。
さらに、任意の二次以下の $g(x)$ について $\dot F=y g(x)$ もこの像に入る。
実際、$x$ に関して二次以下の $P(x,t)$ を
\[
 P(r_i,t)=-h(t,r_i)g(r_i)/f'(r_i)\qquad (1\le i\le3)
\]
で補間し、その $x^2$ 係数を $c(t)$ とする。
$\dot A=P$, $\dot B=y(d(t)+c(t)x)$ は正則な混合項であり、
この変形の一次 Jacobi 条件は
\[
 (f'+h_t)P-hP_t-f(P_x+c x)+h g=f d
\]
となる。左辺は各 $r_i$ で消える。$t$ の微分を含む二項は
二次以下の多項式ベクトル場の括弧であり $t$ の次数は高々2である。
左辺の $x^4$ の係数も $f'P-f(P_x+c x)$ において打ち消し合う。
したがって左辺を $f$ で割れば $t$ の二次以下の多項式 $d$ が得られる。
以上から接空間の次元は少なくとも $6+4+3=13$ である。

最後に一点で階数の逆向きの評価を確認する。
$f=(x-1)(x+2)(x-3)$、
$h=1+2t+3t^2+x(5+7t+11t^2)$ と置く。
Jacobi 式の微分は
\[
 (f'+h_t)\dot A-h\dot A_t-f(\dot A_x+\dot B_y)+h\dot F_y .
\]
$t=0,1,2$、$x,y=0,1,2,3$ の組でこれを評価した行列は
$\mathbb F_{101}$ 上で階数21である
（付属の \path{p1p2/verify_concurrent.py}）。よって標数零で
一般点の接空間は高々 $34-21=13$ 次元である。
台の既約成分の次元は12であり、他の成分との交わりを除けば
この次元差は生成点の非被約性を意味する。
\end{proof}

\begin{theorem}\label{p12:th:double}
$C_{\mathrm{Ric}}$ の生成点 $\eta$ における $\operatorname{Pois}(X)$ の局所環の長さは2である。
\end{theorem}
\begin{proof}
$J$ を30個の斉次二次式の組として用いる。
一次変形 $\delta$ が2次まで延長する条件は、
$J(\delta)\in\operatorname{im} dJ_\pi$ である。実際、
$J(\pi+\epsilon\delta+\epsilon^2\delta_2)$ の $\epsilon^2$ 係数は
$J(\delta)+dJ_\pi(\delta_2)$ である。
前の証明の補間で $g=1$ を取り、
\[
 f=x(x-1)(x-2),\qquad h=1+2t+3t^2+x(4+7t+11t^2)
\]
に特殊化する。ここで $\pi=(0,h,f)$ とし、補間から得る一次変形を明記すると
\[
 \delta=(P,\ y(d+cx),\ y),
\]
\begin{align*}
 P={}&-\tfrac12\bigl[(42t^2+27t+15)x^2
              -(73t^2+47t+26)x+(3t^2+2t+1)\bigr],\\
 c={}&-\tfrac12(42t^2+27t+15),\qquad
 d=\tfrac12(17t^2+19t+10)
\end{align*}
である。$P$ の $x^2$ 係数が $c$ なので、$\delta$ は大域双ベクトルを表す。
\eqref{p12:eq:jacobi} への代入により $dJ_\pi(\delta)=0$ となり、
\[
 J(\delta)=\frac y4\bigl[
 (336t^2+336t-12)x^2+(-417t^2-474t+81)x-123t^2-68t-51\bigr]
\]
を得る。多項式 $R$ の単項式 $m$ の係数を $[m]R$ と書き、線形汎関数
\[
 \ell(R)=\sum_{k,i=0}^2\ell_{ki}[t^kx^iy]R,\qquad
 (\ell_{ki})=
 \begin{pmatrix}
 886&-739&61\\
 -1521&279&-171\\
 756&81&81
 \end{pmatrix}
\]
を定める。線形化の式への代入は
\[
 \ell\circ dJ_\pi=0,\qquad \ell(J(\delta))=-72900\ne0
\]
を与える。従ってこの一次変形には2次障害がある。

階数21の証明書も記す。順序付き単項式列
\[
 \mathcal M=(1,y,y^2,y^3,x,xy,xy^2,x^2,x^2y,x^3)
\]
を用い、$dJ_\pi$ の行を $t^k m$ の順
（$k=0,1,2$、各 $k$ で $m\in\mathcal M$ の順）に並べる。
列は先に $(0,0,m)$ を $\mathcal M$ の順に置き、次に
$(t^kP_j,t^kQ_j,0)$ を並べる。ここで $k=0,1,2$、
$(P_j,Q_j)$ はアフィン表示の節に掲げた八つの平面ベクトル場の係数を
$(v,w)=(x,y)$ としたものである。添字を1から数えると、行・列集合
\begin{align*}
 I={}&(1,2,3,5,6,8,9,10,11,12,15,16,18,19,20,21,22,25,26,28,30),\\
 K={}&(2,3,4,6,7,9,11,13,14,16,17,18,19,21,22,24,25,26,27,30,34)
\end{align*}
で切り出した小行列の行列式は
\[
 \det(dJ_\pi)_{I,K}=-2^{19}3^{13}5^2\cdot11\cdot19\ne0
\]
である。前の証明で構成した13個の独立な接方向と合わせて
$\operatorname{rank}dJ_\pi=21$ が従う。したがって
$\operatorname{rank}[dJ_\pi\mid J(\delta)]=22$ である。
以上の代入と行列式は \path{p1p2/verify_double.py} でも有理数上で検証できる。

補間は $f$ の判別式を逆にすれば係数の有理式である。
固定した座標内の族でこの非零小行列式は開条件であり、
座標変換でその族が $C_{\mathrm{Ric}}$ を稠密に覆う。
従って成分の一般閉点に、2次に延長しない一次変形がある。

一般点の近くで微分が独立な21本の Jacobi 方程式を選ぶ。
その共通零点は滑らかな13次元多様体 $Z$ であり、
$C_{\mathrm{Ric}}$ の被約構造は一般点の近くで $Z$ の滑らかな因子である。
従って $\mathcal O_{Z,\eta}$ は離散付値環で、残余の方程式のイデアルは
一様化元 $s$ のある冪 $(s^m)$ である。他の成分を除く開集合をさらに縮めれば、
そこで同じ等式が成り立つ。定理\ref{p12:th:concurrent-riccati} の接空間次元より $m\ge2$。
$m\ge3$ なら、滑らかな $Z$ 上で任意の接ベクトルを2次まで延長すると
$s^m$ は $\epsilon^3$ の剰余で自動的に零になる。一般点での上の障害に反する。
従って $m=2$ である。
\end{proof}

\subsubsection{三本の相異なる共点直線の上の全解}
\label{p12:sec:concurrentfibers}
\begin{theorem}\label{p12:th:all-concurrent}
平面成分の零曲線が相異なる三本の共点直線の和であるとする。
少なくとも一本の直線の一般点で接続は対数的である。
三本とも対数的なら解は直線束からの引き戻し \eqref{p12:eq:concurrent-family} である。
ちょうど二本が非対数的なら解は $C_{(1,2)(1,0)(0,1)}$ に属する。
ちょうど一本が非対数的なら解は $C_4$ に属する。
\end{theorem}
\begin{proof}
三本を $x=0,1$ と無限遠直線に取る。平面成分を定数倍で
$q=x(x-1)$ に正規化する。射影ベクトル場の二次部分は動径方向にあるから、
水平持ち上げ $\partial_i+\Omega_i$ を表す行列1形式は
\[
 \Omega=(X\,dx+Y\,dy)/q,\quad
 X=A+Mx+Ny+Dxy+Ey^2,\quad
 Y=B+Px+Ry-Dx^2-Exy
\]
と書ける。各係数は $\mathfrak{sl}_2$ の元である。

初めに三本とも非対数的と仮定する。各直線の近傍で法線座標を $s$、
接方向の座標を $z$ として $\Omega=(S\,ds+T\,dz)/s$ と書く。
平坦性から $[S|_{s=0},T|_{s=0}]=T|_{s=0}$ であり、
非対数的なら一般点で $T|_{s=0}$ は非零べき零である。
三直線に沿うこの行列は、零でない有理関数を掛ける自由度の下で順に
\[
 B+Ry,\qquad B+P-D+(R-E)y,\qquad D+Ey'
\]
である。最後は $x'=1/x,\ y'=y/x$ と変換して得る。
各式は線形で、$\mathfrak{sl}_2$ のべき零元の射影化は滑らかな二次曲線であるから、
それぞれ一定のべき零方向 $e_i$ を張る。

この三方向は互いに異なる。実際、共点を通らない一般の直線 $y=ax+b$ に制限すれば、
自明な階数2束の行列1形式で三つの単純極だけをもつ。その留数 $S_i$ は
$\sum_i S_i=0$ を満たし、接方向項の寄与は $e_i$ の倍数なので
$[S_i,e_i]=e_i$ となる。トレース零より $S_i$ は $\ker e_i$ 上に $1/2$ で作用し、
固有値は $\pm1/2$ である。二方向が一致すれば第三の留数は同じ直線上で
$-1$ として作用するので矛盾する。

ところが、三つの相異なるべき零方向は滑らかな射影二次曲線上の三点なので
線形独立である。$R,R-E,E$ はそれぞれこの三方向上にあって
$R-(R-E)-E=0$ だから $R=E=0$ である。
$Y=Y_0(x)=B+Px-Dx^2$ とおくと、平坦性の $y$ の係数は $[N+Dx,Y_0(x)]=0$ となる。
$Y_0(0),Y_0(1)$ は非零べき零で、その $\mathfrak{sl}_2$ における中心化代数は
それぞれの張る直線である。従って $N,N+D,D$ はそれぞれ $e_0,e_1,e_\infty$
の方向にあり、同じ線形独立性から $D=0$ となる。
これは無限遠直線で非対数的であることに反する。

対数的な直線を無限遠に選び直すと $D=E=0$ となる。
$X=A+Mx+Ny,\ Y=B+Px+Ry$ を平坦性の式に代入して
\begin{align*}
 [M,P]&=P+N, & [M,R]+[N,P]&=2R,&[N,R]&=0,\\
 N-2B+[M,B]+[A,P]&=0, & R+[N,B]+[A,R]&=0,&[A,B]&=-B
\end{align*}
を得る。$x=0,1$ の接方向の係数は対数的なら零であり、いずれの場合にも
$B+Ry$ と $B+P+Ry$ は全ての $y$ に対してべき零である。

$R\ne0$ なら $R,B,P$ は同じべき零方向 $e$ を持ち、
両直線とも非対数的である。$[N,R]=0$ から $N$ も $e$ の倍数となる。
$[M,R]=2R,\ [A,R]=-R$ から $M=2h+m e,\ A=-h+a e$、ついで $N=P$ を得る。
ここで $[h,e]=e$ とし、$h,e$ は前節と同じ標準形に取る。
よって $P=p e,\ R=r e,\ B=b e$ と書くと
\[
  \Omega=h\,dq/q+e\,dg/q,\qquad
  g=mx^2/2+p xy+r y^2/2+ax+by.
\]
これは有理第一積分 $q/t+g$ の $(1,2)$ 型の族である。

$R=0$ とする。べき零元 $B,B+P$ が一次独立なら $[N,B]=[N,P]=0$ から $N=0$ である。
このとき $P$ は独立な二つのべき零元の差なので $\operatorname{tr}(P^2)\ne0$ だが、
$[M,P]=P$ に反する。従って $B=b e,\ P=p e$ と置ける。
$p=0$ なら $N=0$ である。$b\ne0$ なら $[M,B]=2B,\ [A,B]=-B$
から再び上の $(1,2)$ 型を得る。$p=b=0$ なら三本とも対数的で、
$Y=0,\ X=A+Mx$ が全解である。

$p\ne0$ なら $[N,P]=0$ なのである $m,u,v,a\in\mathbb C$ を用いて
\[
 M=m h+u e,\qquad N=(m-1)p e,\qquad A=a h+v e
\]
と書ける。実際、最初の式は $[M,P]=P+N$ を比較し、$A$ の式は
$N-2B+[M,B]+[A,P]=0$ を比較して得る。
残る条件は $b(a+1)=0,\ p(a+m-1)=b(2-m)$ である。
$b(b+p)\ne0$ なら両直線とも非対数的で、$a=-1,m=2$ となり上の $(1,2)$ 型である。
残る場合には $0$ と $1$ を交換すれば $b=0$ とでき、このとき $a=1-m$ である。
$k=m-1$ と置くと、接続の係数は
\[
 h\left(k\frac{dx}{x}+\frac{dx}{x-1}\right)
 +e\left(\frac{ux+kpy+v}{x(x-1)}\,dx+\frac{p}{x-1}\,dy\right).
\]
$k\ne0,-1$ とし、$\ell(x)=u x/(k+1)+v/k$ と置く。
$h=t\partial_t,\ e=t^2\partial_t$ で水平方程式を書けば、
$H=(x-1)/t+p y+\ell(x)$ は $dH+kH\,dx/x=0$ を満たす。
$H$ を $F/(U L_\infty)$ と斉次表示すると $U=t$、$F$ は $(1,1)$ 型で、
$x$ の零直線と $L_\infty$ を合わせて \eqref{p12:eq:four} の四因子を得る。
分母を払った1形式の $dt$ 係数は $q$ の非零定数倍なので、
この表示は大域双ベクトル自身の定数倍までの一致を与える。
$k=0,-1$ でも、$p,u,v$ を固定して $k$ を動かした元の接続は $k$ の多項式族であり、
解は閉包 $C_4$ に入る。
\end{proof}


\subsection{二重直線と単純直線}
\subsubsection{二重直線と単純直線のうち単純直線で対数的な場合}
\begin{proposition}\label{p12:th:double-good}
平面成分の零因子を $2L+M_0$ とし、$L\ne M_0$ とする。
単純直線 $M_0$ の一般点で接続が対数的なら、解は
$C_{(1,2)(1,0)(0,1)}, C_{\mathrm{Ric}}, C_4, C_{\mathrm{tor}}$
のいずれかに属する。
\end{proposition}
\begin{proof}
$M_0$ を無限遠直線、$L=(x=0)$ として平面成分を $q=x^2$ に正規化する。
$M_0$ における対数性より、水平接続を定める行列1形式は
\[
 \Omega=((A+Bx+Cy)\,dx+(D+Ex+Gy)\,dy)/x^2
\]
である。各係数は $\mathfrak{sl}_2$ の元であり、平坦性は
\begin{align*}
 [B,E]&=E+C,&[B,G]+[C,E]&=2G,&[C,G]&=0,\\
 [B,D]+[A,E]&=2D,&[C,D]+[A,G]&=0,&[A,D]&=0
\end{align*}
に等しい。とくに全ての $y$ で $[A+Cy,D+Gy]=0$ である。

$D,G$ が一次独立なら、$[A,D]=[C,G]=0$ および中心化代数の一次元性から
$A=aD,C=cG$ と書ける。混合係数の消滅と $[D,G]\ne0$ より $a=c$ となる。
$y'=y+ax$ と変換すると、新しい表示の $A,C$ は零になり、$D,G$ は変わらない。
ところがこの表示で $[B,D]=2D,\ [B,G]=2G$ を満たすのは不可能である。
実際 $\mathfrak{sl}_2$ の随伴作用の非零固有値に対する固有空間は高々一次元である。

$D,G$ は一次従属で少なくとも一方は非零とする。
$D=dK,\ G=gK$ と置くと、先の交換関係から
$A=aK,\ C=cK$ と書ける。従って
\[
 g[B,K]+c[K,E]=2gK,\qquad d[B,K]+a[K,E]=2dK.
\]
$(d,g),(a,c)$ が一次独立なら $[B,K]=2K,\ [K,E]=0$ である。
よって $K$ はべき零、$E$ は $K$ の定数倍であり、$[B,E]=E+C$ から $C=E$ を得る。
$[h,e]=e,\ K=e$ の標準表示を取ると $B=2h+b e$ となる。
すると $h$ 成分は $dq/q$ で、$e$ 成分の分子は $C=E$ によって閉じたアフィン線形1形式である。
従ってある二次以下の多項式 $g_0$ について
$\Omega=h\,dq/q+e\,dg_0/q$ となり、第一積分 $q/t+g_0$ をもつ $(1,2)$ 型の族に入る。
二組が比例すれば再び $y'=y+\lambda x$ により $A=C=0$ とできる。
$[B,K]=2K,\ [B,E]=E$ が同時に成り立つことから $E=0$ であり、同じ形の解を得る。

以下 $D=G=0$ とする。$E=0$ なら $C=0$ で、$\Omega=(A+Bx)dx/x^2$ は
直線束からの引き戻しである。射影直線上の binary cubic の重根を動かして
\eqref{p12:eq:concurrent-family} からの特殊化と書ける。
$E\ne0$ なら $[A,E]=[C,E]=0$ から $A=aE,C=cE$ と書け、
$[B,E]=(1+c)E$ となる。

$c\ne-1$ なら $E$ はべき零で、$E=e,\ B=(1+c)h+b e$ と置ける。
接続は
\[
 (1+c)h\,\frac{dx}{x}+
 e\left(\frac{a+bx+cy}{x^2}\,dx+\frac{dy}{x}\right).
\]
さらに $c\ne0$ とする。$h=t\partial_t,\ e=t^2\partial_t$ とすれば
\[
 W=\frac1t+\frac{y+a/c}{x}+\frac{b}{1+c}
\]
は水平方向で $dW+(1+c)W\,dx/x=0$ を満たす。
$W=F/(U L)$ と斉次表示でき、アフィン式は $U=t$, $L=x$,
$F=x+t(y+a/c+bx/(1+c))$ である。
分母を払うと $(1,1)$ 型の $F$、$U$、$L$、無限遠直線を用いた
四因子形式 \eqref{p12:eq:four} が得られる。実際その $dt$ 係数は
$x^2$ の非零定数倍である。$c=0$ でも $a,b$ を固定して $c$ を動かせば、
元の行列表現が $c$ に多項式的に依存するので $C_4$ への所属が従う。

最後に $c=-1$ なら $B=bE$ であり、
\[
 \Omega=E\left(b\,d\log x+d\frac{y-a}{x}\right)
\]
となる。任意の $E$ についてこれがトーラス成分に属することは、三角形からの
次の明示的特殊化で分かる。$\delta\ne0$ とし、
\[
 q_\delta=x(x+\delta(y-a)),\qquad
 \Omega_\delta=E\left(b\,d\log x+
 \delta^{-1}d\log\frac{x+\delta(y-a)}{x}\right)
\]
とする。極の三直線は非共点である。平面のトーラス座標を $x,x+\delta(y-a)$
と取れば、この接続は同じ $E$ の定数倍に値を取る不変対数形式である。
$E$ が半単純なら両因子のトーラス作用から生じ、一般の $E$ も半単純元の極限である。
$q_\delta\Omega_\delta$ の分母を払うと $\delta^{-1}$ は消え、
$\delta\to0$ で $q_\delta\to q$, $q_\delta\Omega_\delta\to q\Omega$
となるので、大域双ベクトルの特殊化が得られる。
\end{proof}

\subsubsection{二重直線と単純直線のうち単純直線で非対数的な場合}
\begin{proposition}\label{p12:prop:doublebad}
射影三次曲線が $2L+M$（$L\ne M$ は直線）であり、単直線 $M$ の一般点で
接続が非対数的であるとする。この解は $C_4$ または
$C_{(1,2)(1,0)(0,1)}$ に属する。
\end{proposition}
\begin{proof}
$L$ を無限遠、$M$ を $x=0$ に取り、平面成分を $q=x$ とする。
接続行列の分子の二次部分は射影ベクトル場の動径成分から生じ、
行列を
\[
 \Omega=\frac{(a+bx+cy+mxy+ny^2)\,dx+
                (d+ex+fy-mx^2-nxy)\,dy}{x}
\]
と書ける。各係数は $\mathfrak{sl}_2$ の元である。
$x=0$ が非対数的なので $d+fy$ は一般点で非零であり、常にべき零である。
従って $d=DJ$, $f=FJ$ となる一定のべき零方向
$J=E_{12}$ があり $(D,F)\ne(0,0)$ である。平坦性の二位の極から
$[a+cy+ny^2,J]=J$ を得る。従って $c=CJ$, $n=NJ$ であり、
$J$ を固定する定数共役により $a=H=\operatorname{diag}(1/2,-1/2)$ とできる。

二次成分の消滅を示す。$S=H+(Cy+Ny^2)J$ とおくと
\[
 \Omega_x=S/x+b+my,\qquad
 \Omega_y=(D+Fy)J/x+e-mx-NJy.
\]
平坦性 $\partial_x\Omega_y-\partial_y\Omega_x+
[\Omega_x,\Omega_y]=0$ で $x$ の係数および定数項を比較すると
\[
 [b,m]=0,\qquad -2m-[H,m]+[b,e]=0,\qquad
 -C[J,m]-N[b,J]+[m,e]=0 .
\]
$m\ne0$ なら $b=\lambda m$ と書け、これらは
$[H+\lambda e,m]=-2m$ および $[m,e+(C-N\lambda)J]=0$ を与える。
後者から $e=-(C-N\lambda)J+\mu m$ である。
従ってある定数 $s$ について $[H+sJ,m]=-2m$ を得るが、
$\operatorname{ad}(H+sJ)$ の固有値は $0,\pm1$ なので矛盾する。
ゆえに $m=0$ である。上記の定数項と $x^{-1}$ の係数の
$y$ に関する各係数は
\begin{align*}
 [b,e]=0,\qquad N[b,J]=0,\qquad [H,e]+D[b,J]=CJ,\\
 C[J,e]+F[b,J]=3N J,\qquad N[J,e]=0
\end{align*}
となる。$N\ne0$ なら左辺の二つの括弧がともに零となり $3NJ=0$ に反する。
従って $n=0$ であり、すべての混合ベクトル場は $L$ に接する。

次に全行列が同時に上三角であることを示す。$b,e$ は可換なので
$b=\lambda Z,e=\mu Z$ と書ける。$Z$ の下三角成分が非零とすると
$[H,e]+D[b,J]=CJ$ の下三角成分から $\mu=0$ が従う
（任意の $Z$ で $[Z,J]$ の下三角成分は零）。
$[Z,J]$ の対角成分は非零なので同じ等式から $D\lambda=0$ かつ $C=0$、
ついで $C[J,e]+F[b,J]=0$ より $F\lambda=0$ である。
$(D,F)\ne(0,0)$ なので $b=e=0$ を得る。従っていずれにせよ全係数は上三角である。

従ってあるアフィン一次式 $\ell$、次数1以下の多項式を係数とする
平面1形式 $\xi$ を使い
\[
 \Omega=H\,(d\log x+d\ell)+J\,\xi/x
\]
と書ける。水平な射影座標を $(t,1)^{\mathsf T}$ とすれば、
$s=xt$ で水平方程式は $ds+s\,d\ell+\xi=0$ となり、
$d\xi+d\ell\wedge\xi=0$ である。
$d\ell=0$ なら $\xi=dP$ をみたす次数2以下の多項式 $P$ がある。
このとき $xt+P$ の斉次化が $(1,2)$ 型で、定理\ref{p12:th:threecomponents} 第1族
（$U,L$ を無限遠に選ぶ）の水平方程式を得る。

$d\ell\ne0$ なら次数1以下の $P$ で $\xi=dP+P\,d\ell$ と書ける。
実際に $d_\ell=d+d\ell\wedge$、$d\ell(Z)=1$ となる定数ベクトル場
$Z$ をとると $d_\ell\iota_Z+\iota_Zd_\ell=1+\mathcal L_Z$ である。
定数方向微分 $\mathcal L_Z$ は多項式係数上局所べき零なので、
$P=(1+\mathcal L_Z)^{-1}\iota_Z\xi$ は同じ次数の上限を満たす。
$F=xt+P$ の斉次化は $(1,1)$ 型で、水平方程式は $dF+F\,d\ell=0$ となる。
$\ell=L'/L$ とし、$U$ をファイバーの無限遠とすると
これは四因子対数族の特殊化である。すなわち
\[
 FUL(L+\epsilon L')\left(
 d\log\frac{F}{UL}+\epsilon^{-1}d\log\frac{L+\epsilon L'}{L}\right)
\]
は $\epsilon\ne0$ で \eqref{p12:eq:four} に属し、
$U=L=1$ のチャートでの極限は $dF+F\,d\ell$ である。
いずれの場合も $dt$ 成分が $x$ であるから、同じ平面成分をもつ
双ベクトルとしても一致する。
\end{proof}

\subsection{三重直線}
\begin{proposition}\label{p12:prop:triple-cubic}
平面成分の零因子が三重直線の場合、対応する点は
$C_{\mathrm{Ric}}$、$C_{\mathrm{line\text{-}conic}}$、
$(1,0)+(1,1)+(0,1)+(0,1)$ 型の対数形式の閉包の和集合に入る。
\end{proposition}
\begin{proof}
三重線を無限遠直線とし、スカラー倍で $\sigma=\partial_x\wedge\partial_y$ とする。
射影ベクトル場の二次の部分が一次式と Euler ベクトル場の積であることから
$\Omega=P\,dx+Q\,dy$ の斉次二次部分は
$L(y\,dx-x\,dy)$、$L=xH+yI$ となる。行列はすべてトレース零とする。
$P,Q$ の一次部分を $P_1,Q_1$ とし、$M=xP_1+yQ_1$ とおく。
曲率の三次部分は $-[M,L]$ である。$L\ne0$ と仮定する。

まず $H,I$ が独立なら、$\mathbb C(x,y)$ 上での非零トレース零行列の中心化環は
$\mathfrak{sl}_2$ 内で1次元なので、$M$ は $L$ の有理関数倍である。
$L$ の成分に公約因子がないので $M=(a x+b y)L$ と書ける。
平行移動 $x\mapsto x+r,\ y\mapsto y+s$ によって $M$ は
$(s x-r y)L$ だけ変化する。実際、二次部分に由来する一次部分の増分は
$(rH+sI)(y\,dx-x\,dy)+L(s\,dx-r\,dy)$ で、Euler ベクトル場に代入すればよい。
従って $M=0$ にでき、$\Omega=(L+C)(y\,dx-x\,dy)+A\,dx+D\,dy$ となる。
曲率二次項より $[xA+yD,L]=0$、したがって $xA+yD=kL$（$k$ は定数）。
一次項は $-3L+k[C,L]=0$ であり、$k\ne0$ かつ $H,I$ が
$\operatorname{ad}C$ の同じ非零固有値に属することになる。
この固有空間は高々1次元なので矛盾する。

次に $H,I$ が従属で $L\ne0$ なら、線形座標変換と定数の取り替えで $L=xH$ とできる。
$[M,L]=0$ により $M$ は二次式と $H$ の積であり、前と同じ平行移動により
$M=g y^2 H$ とできる。従って
\[
 \Omega=(xH+C)(y\,dx-x\,dy)+gHy\,dy+A\,dx+D\,dy .
\]
曲率二次項の $x^2,xy$ の係数は $-[A,H],[H,D]$ だから $A,D$ は $H$ の倍数である。
定数項 $-2C+[A,D]=0$ より $C=0$ となり、一次項 $-3xH$ が消えない。
よってすべての場合に接続の二次部分は零である。

今や $P=A+P_1$, $Q=D+Q_1$ であり、曲率二次項は $[P_1,Q_1]=0$ を与える。
$P_1$ の二つの行列係数が独立なら、先の中心化環の議論から $Q_1=kP_1$ である。
底の線形座標変換で一次部分を $P'_1\,dx$ にでき、その二つの係数はなお独立である。
平坦性より $[P'_1,D']=0$ なので $D'=0$、さらに $\partial_yP'_1=0$ となり矛盾する。
$Q_1$ についても同じである。従って一次部分が非零なら、その全行列係数は共通の
行列 $H$ の倍数である。$P_1=lH,Q_1=mH$ で一次式 $l,m$ が独立なら
一次の曲率条件より $A,D$ も $H$ の倍数となる。従属な場合は座標変換により
\[
 \Omega=(A+(bx+cy)H)dx+D\,dy,\qquad
 D=dH,\qquad d[A,H]=cH
\]
と書ける。一次部分が零の場合は平坦性が $[A,D]=0$ を意味する。

全係数が可換な場合には $\Omega=H\,dK$, $\deg K\le2$ である。
$1+\epsilon K$ をアフィン二次式、$L$ を無限遠直線として
$\mathrm{line\text{-}conic}$ 構成を適用すると、
ファイバーの多項式を $1/\epsilon$ 倍することでこの形を極限に得る。
ここで $K$ が退化していても一般の二次式からの特殊化で従う。
係数が可換でなく $c=0$ の場合は $D=0$ で、
$(A+b xH)dx$ は相異なる三根を持つ $f_\epsilon\to1$ による
$(A+b xH)dx/f_\epsilon$ の極限だから $C_{\mathrm{Ric}}$ に属する。

最後に $c\ne0$ なら $\alpha=c/d\ne0$ とおく。
$[A,H]=\alpha H$ によりファイバーの座標変換で
$H=E_{12}$, $A=\alpha\,\mathrm{diag}(1/2,-1/2)$ とできる。
水平式は $dt+\alpha t\,dx+(bx+cy)\,dx+d\,dy=0$ である。
$T=t+d y+(b/\alpha)x-b/\alpha^2$ とおくとこの式は $dT+\alpha T\,dx=0$ となる。
これは $(1+\epsilon x)dT+\alpha Tdx$ の極限である。
後者は $d\log T+(\alpha/\epsilon)d\log(1+\epsilon x)$ の分母を払った形式で、
$T$ の分子は $(1,1)$ 型、分母はファイバーと平面の無限遠切断の積である。
もう一つの直線 $1+\epsilon x=0$ と合わせて主張の四因子型となる。
\end{proof}


\subsection{既約成分の網羅性}\label{p12:sec:conclusion}
\begin{proof}[定理\ref{p12:th:complete}の証明]
平面三次式 $C$ を固定する。$C=0$ のときは命題\ref{p12:prop:zero-cubic}、
被約既約の場合は定理\ref{p12:th:irreducible-cubic} がファイバーを尽くす。
$C=LQ$ で $Q$ が滑らかな二次曲線なら定理\ref{p12:th:wholeconic}、
相異なる三直線の積なら定理\ref{p12:th:triangle}, \ref{p12:th:all-concurrent} が適用できる。
$C=L^2M$ で $L\ne M$ なら命題\ref{p12:th:double-good}, \ref{p12:prop:doublebad} が単純成分に沿う二つの場合を尽くし、
$C=L^3$ は命題\ref{p12:prop:triple-cubic} で尽きる。
これらは複素数上の全三次式を覆う。各場合、すべての点は上表の九つの既約閉集合の
いずれかに属する。他方、各閉集合が実際に成分であることは、 $C_{\mathrm{Ric}}$ 以外の八つについては
像の次元と接空間による証明から、$C_{\mathrm{Ric}}$ については
定理\ref{p12:th:concurrent-riccati} から従う。
$C_{\mathrm{Ric}}$ は前の二つの12次元成分と異なる。実際、その一般点で平面三次式は
共点三直線を定め、他では非共点三直線または直線と滑らかな二次曲線を定める。
台がそのような共点三直線またはその特殊化となる条件は閉である
（中心を動かす二変数三次式の射影パラメータ空間の像）。
他の成分同士の相異性は命題\ref{p12:prop:distinct} で示した。
前の八成分の各々には全 Poisson 概型の滑らかな点があり、従って生成点での長さは1である。
$C_{\mathrm{Ric}}$ の長さは定理\ref{p12:th:double} による。
\end{proof}

この分類は固定した $X$ 上のアフィン概型についてのものである。
異なる成分の全交叉イデアル、埋め込まれた素因子の有無、
および一般の自己同型軌道の分類は未解決である。
各成分の階数証明書、固定因子の制限、四因子型の Smith 不変因子、
共点三直線型の2次障害は同梱のコードで再現できる。
